\documentclass{article}

% if you need to pass options to natbib, use, e.g.:
%     \PassOptionsToPackage{numbers, compress}{natbib}
% before loading neurips_2025

% The authors should use one of these tracks.
% Before accepting by the NeurIPS conference, select one of the options below.
% 0. "default" for submission
 % \usepackage{neurips_2025}
% the "default" option is equal to the "main" option, which is used for the Main Track with double-blind reviewing.
% 1. "main" option is used for the Main Track
%  \usepackage[main]{neurips_2025}
% 2. "position" option is used for the Position Paper Track
%  \usepackage[position]{neurips_2025}
% 3. "dandb" option is used for the Datasets & Benchmarks Track
 % \usepackage[dandb]{neurips_2025}
% 4. "creativeai" option is used for the Creative AI Track
%  \usepackage[creativeai]{neurips_2025}
% 5. "sglblindworkshop" option is used for the Workshop with single-blind reviewing
 % \usepackage[sglblindworkshop]{neurips_2025}
% 6. "dblblindworkshop" option is used for the Workshop with double-blind reviewing
%  \usepackage[dblblindworkshop]{neurips_2025}

% After being accepted, the authors should add "final" behind the track to compile a camera-ready version.
% 1. Main Track
 \usepackage[main, final]{neurips_2025}
% 2. Position Paper Track
%  \usepackage[position, final]{neurips_2025}
% 3. Datasets & Benchmarks Track
 % \usepackage[dandb, final]{neurips_2025}
% 4. Creative AI Track
%  \usepackage[creativeai, final]{neurips_2025}
% 5. Workshop with single-blind reviewing
%  \usepackage[sglblindworkshop, final]{neurips_2025}
% 6. Workshop with double-blind reviewing
%  \usepackage[dblblindworkshop, final]{neurips_2025}
% Note. For the workshop paper template, both \title{} and \workshoptitle{} are required, with the former indicating the paper title shown in the title and the latter indicating the workshop title displayed in the footnote.
% For workshops (5., 6.), the authors should add the name of the workshop, "\workshoptitle" command is used to set the workshop title.
% \workshoptitle{WORKSHOP TITLE}

% "preprint" option is used for arXiv or other preprint submissions
 % \usepackage[preprint]{neurips_2025}

% to avoid loading the natbib package, add option nonatbib:
   % \usepackage[nonatbib]{neurips_2025}

\usepackage[utf8]{inputenc} % allow utf-8 input
\usepackage[T1]{fontenc}    % use 8-bit T1 fonts
\usepackage{hyperref}       % hyperlinks
\usepackage{url}            % simple URL typesetting
\usepackage{booktabs}       % professional-quality tables
\usepackage{amsfonts}       % blackboard math symbols
\usepackage{nicefrac}       % compact symbols for 1/2, etc.
\usepackage{microtype}      % microtypography
\usepackage{xcolor}         % colors

\usepackage{graphicx}      % figures
\usepackage{subcaption}    % subfigure env
\usepackage{float}         % H option for figure
\usepackage{wrapfig}
\usepackage{caption}       % captions
\usepackage{natbib} 

\usepackage{amsmath, amssymb, amsthm}
\usepackage{thmtools}
\usepackage{bbm}  % for \mathbbm{1}

\usepackage{algorithm}      % algorithm
\usepackage{algorithmic}    % algorithm

\usepackage{caption}
\usepackage{multirow}
\usepackage{colortbl}    % \cellcolor
\newcommand{\cc}[1]{\cellcolor{gray!#1}}

\newcommand{\BarStar}{%
  \rule{0.8cm}{0.25cm}%
  \llap{%
    \hspace{0.4cm}%
    \raisebox{0.02cm}{\textcolor{white}{*}}%
  }%
}

\newcounter{dummylabelcounter}

% for general theorem
\newtheorem{theorem}{Theorem}

% theorem where we want to replicate
% \newtheorem*{reptheorem}{Theorem}
\newtheorem*{reptheorem}{\textnormal{\textbf{Theorem \ref{thm:safe}}}}

\newcommand{\diff}{\mathop{}\!\mathrm{d}} % diff 


\usepackage{pifont}     % for \ding{55} (xmark)
\newcommand{\cmark}{\checkmark}          % check-mark
\newcommand{\xmark}{\ding{55}}           % x-mark

% --- customized tcolorbox env. ---
\usepackage[most]{tcolorbox}
\usepackage{xcolor}

\definecolor{DefinitionBackColor}{gray}{0.96}
\definecolor{DefinitionBarColor}{gray}{0.5}

\newtcolorbox{mydefinitionbox}{
    enhanced,
    breakable,
    colback=DefinitionBackColor,
    colframe=DefinitionBarColor,
    boxrule=0pt, leftrule=3pt,
    rightrule=0pt, toprule=0pt, bottomrule=0pt,
    arc=0mm,
    boxsep=0pt, top=4pt, bottom=4pt, left=8pt, right=4pt,
    before skip=\medskipamount, after skip=\medskipamount,
    fontupper=\normalsize,
    fonttitle=\bfseries,
}

% Note. For the workshop paper template, both \title{} and \workshoptitle{} are required, with the former indicating the paper title shown in the title and the latter indicating the workshop title displayed in the footnote. 
\title{Training-Free Safe Denoisers \\for Safe Use of Diffusion Models}


% The \author macro works with any number of authors. There are two commands
% used to separate the names and addresses of multiple authors: \And and \AND.
%
% Using \And between authors leaves it to LaTeX to determine where to break the
% lines. Using \AND forces a line break at that point. So, if LaTeX puts 3 of 4
% authors names on the first line, and the last on the second line, try using
% \AND instead of \And before the third author name.


\author{%
  Mingyu Kim\thanks{Equal contribution}\ $^{1}$ \ Dongjun Kim$^{*2}$ \ Amman Yusuf$^1$ \ Stefano Ermon$^2$ \ Mijung Park$^1$\\
  $^1$CS, UBC \ \ $^2$CS, Standford\\
  \texttt{mgyu.kim@ubc.ca}, \ \ 
  \texttt{dongjun@stanford.edu} \\
  \texttt{ammany01@cs.ubc.ca},  \ \
  \texttt{ermon@cs.stanford.edu},  \ \
  \texttt{mijungp@cs.ubc.ca} \\
}

% math expression
\input{math_commands.tex}

\begin{document}


\maketitle


\begin{abstract}
  There is growing concern over the safety of powerful diffusion models, as they are often misused to produce inappropriate, not-safe-for-work content or generate copyrighted material or data of individuals who wish to be forgotten. 
  Many existing methods tackle these issues by heavily relying on text-based negative prompts or retraining the model to eliminate certain features or samples. 
  In this paper, we take a radically different approach, directly modifying the sampling trajectory by leveraging a negation set (e.g., unsafe images, copyrighted data, or private data) to avoid specific regions of data distribution, without needing to retrain or fine-tune the model. 
  We formally derive the relationship between the expected denoised samples that are safe and those that are unsafe, leading to our \textit{safe} denoiser, which ensures its final samples are away from the area to be negated.
  We achieve state-of-the-art safety in large-scale datasets such as the CoPro dataset
  while enabling significantly more cost-effective sampling than existing methodologies.
%  Inspired by the derivation, we develop a practical algorithm that successfully produces high-quality samples while avoiding negation areas of the data distribution in text-conditional, %class-conditional, and unconditional image generation scenarios. These results hint at the great potential of our training-free safe denoiser for using DMs more safely.
\end{abstract}

\vspace{-0.4cm}
\textcolor{red}{Warning: This paper contains disturbing content such as violent and sexually explicit images.}
%
\begin{figure}[H]
    \centering
    \vskip -0.1in
        \begin{subfigure}{0.9\textwidth}
            \includegraphics[width=\textwidth]{Figures/Introduction/thumbnail_v3.pdf}
            \vskip -0.05in
            \caption{Nudity (Top)}
            \label{fig:thumbnail_1}
        \end{subfigure}    
        \begin{subfigure}{0.7\textwidth}
                \includegraphics[width=\textwidth]{Figures/Introduction/thumbnail_ann_v2.pdf}
                \vskip -0.05in
            \caption{Data Memorization}
            \label{fig:thumbnail_2}
        \end{subfigure}
        \vskip -0.05in
    \caption{Our method \textit{Safe Denoiser} against existing methods. (a) Our method, incorporated with SAFREE~\cite{yoon2024safree} and SLD~\cite{schramowski2023safe}, does not generate inappropriate images. (b) Our method mitigates the memorization issue by negating the real image, resulting in a novel image with features similar to those in the real image in hair colors or outfits.}
    \label{fig:thumbnail}
\end{figure}
\section{Introduction}
\label{sec_introduction}

\begin{figure*}[t]{
    \centering
        \begin{subfigure}{0.4\textwidth}
            \includegraphics[width=\textwidth]{Figures/Method/algorithm.pdf}
            \caption{Schematic overview}
        \end{subfigure}    
        \begin{subfigure}{0.45\textwidth}
            \includegraphics[width=\textwidth]{Figures/Method/moons_trial_0_for_publication_v2.pdf}
            \caption{2-dimensional trajectory}
        \end{subfigure}
    \caption{An overview of the safe denoiser. (a) The safe denoiser $\mathbb{E}_{\text{safe}}$ negates the direction of the unsafe denoiser $\mathbb{E}_{\text{unsafe}}$ from the data denoiser $\mathbb{E}_{\text{data}}$. 
    (b) Trajectories from data denoiser and safe denoiser, starting from the same initial point far from the data distribution, reveal distinct paths: while the sample path from the data denoiser falls into the unsafe region, the trajectory from the safe denoiser successfully avoids it.}
    \label{fig:schematic}
}
\end{figure*}

Diffusion models (DMs) have become leading generative models, excelling in generation tasks like text-to-image~\cite{rombach2022high}, audio~\cite{kong2021diffwave}, video~\cite{bartal2024lumiere}, and protein design~\cite{watson2023denovo}, thanks to their flexible and controllable sampling~\cite{dhariwal2021diffusion, ho2021classifier}. However, growing concerns over unsafe content — such as not-safe-for-work(NSFW) imagery (\figref{thumbnail_1}), copyright violations, and potential misuse — highlight the need for safety. The key challenge is mitigating these risks without compromising model utility or creativity.

Mainstream mitigation strategies for issues like NSFW content or unwanted concept removal rely on text-based guidance~\cite{gandikota2023erasing, ban2024understanding, yoon2024safree} or fine-tuning for unlearning~\cite{gandikota2023erasing, gandikota2023unified, gong2024reliable, kim2024race}. Text-based methods require iterative, expert-crafted negative prompts~\citep{schramowski2023safe}, which may not generalize well, while fine-tuning is resource-intensive and risks catastrophic forgetting or degraded performance on desired tasks.

Other significant safety concerns involve the DMs' capacity to reproduce copyrighted content and their generation of data pertaining to individuals in \figref{thumbnail_2} who wish to be excluded. These issues are often linked to the models' remarkable ability to memorize training data \cite{carlini2023extracting}. While techniques like differentially private training \cite{dockhorn2023differentially, liu2024differentially} can formally limit memorization by adding noise during the training process, they often result in a noticeable degradation in generation quality, which can be particularly prohibitive for applications demanding high-fidelity outputs.

We propose a \textit{safe denoiser} (defined in Definition~\ref{thm:def}) that modifies sampling trajectories such that the resulting samples are drawn from a safe distribution 
(shown in \figref{schematic}). The intuion comes from our \thmref{safe}, 
% which relates safe and unsafe denoised samples, 
where the safe denoiser steers generation away from unsafe regions, ensuring theoretical safety. We develop a practical algorithm (\algoref{safer}) based off of our theorem, which can be used standalone or combined with negative prompting to enhance safety in text-to-image generation. 
% Empirically, 
Our method achieves state-of-the-art performance on concept erasing, class removal, and unconditional image generation tasks.
\section{Preliminary}
\label{preliminary}



%
%
DMs generate samples through iterative 
decoding starting from random noise to data. This iterative process is a reverse of the forward 
data corruption 
%\mj{descruction -> isn't corruption better?} 
process, 
% given by  
$\mathbf{x}_{t}=\alpha_{t}\mathbf{x}+\sigma_{t}\bm{\epsilon}$, where $\mathbf{x} \sim p_{\text{data}}(\mathbf{x})$ 
% follows the data distribution $p_{\text{data}}(\mathbf{x})$ and 
$\bm{\epsilon} \sim \mathcal{N}(0,I)$ 
% follows the noise prior distribution $\mathcal{N}(0,I)$, 
which results in a perturbation kernel: $q_{t}(\mathbf{x}_{t}\vert\mathbf{x})=\mathcal{N}(\mathbf{x}_{t};\alpha_{t}\mathbf{x},\sigma_{t}^{2}I)$. The specific choice of the coefficients $\alpha_{t}$ and $\sigma_{t}$ determines a different variant of DMs: 
% In other words, depending on these parameters, the model may be referred to as 
popular examples include 
Denoising Diffusion Probabilistic Models (DDPM) \cite{ho2020denoising}, Elucidating Diffusion Models (EDM) \cite{karras2022elucidating}, or Flow Matching \cite{lipman2022flow}. 
%
Regardless of whether the model is trained with noise-prediction~\cite{ho2020denoising}, data-prediction~\cite{karras2022elucidating}, or velocity-prediction~\cite{salimans2022progressive,lipman2022flow}, these approaches are fundamentally equivalent~\cite{kingma2021variational,kim2021soft}. This paper adopts the data-prediction framework due to its most intuitive interpretation. In data-prediction, the model approximates the \textit{denoiser} function, defined by
$\mathbb{E}_{\text{data}}[\mathbf{x}\vert\mathbf{x}_{t}]:=\int \mathbf{x} \frac{p_{\text{data}}(\mathbf{x})q_{t}(\mathbf{x}_{t}\vert\mathbf{x})}{p_{\text{data},t}(\mathbf{x}_{t})}\diff\mathbf{x}\approx\frac{1}{\alpha_{t}}(\mathbf{x}_{t}-\sigma_{t}\bm{\epsilon}_{\bm{\theta}})$,
where $p_{\text{data},t}(\mathbf{x}_{t})$ is a marginal distribution of diffusion process at $t$, and $\bm{\epsilon}_{\bm{\theta}}$ is the noise-prediction.

DMs can be guided to produce samples~\cite{dhariwal2021diffusion,kim2022refining} that adhere more closely to a desired condition denoted by $\mathbf{c}$. A common approach in modern DMs is \textit{classifier-free guidance} (CFG)~\cite{ho2021classifier}. The model is trained to learn both the unconditional denoiser $\mathbb{E}_{\text{data}}[\mathbf{x}\vert\mathbf{x}_{t}]$ and the conitional denoiser $\mathbb{E}_{\text{data}}[\mathbf{x}\vert\mathbf{x}_{t},\mathbf{c}]$. The CFG modifies the sampling trajectory by
\begin{align*}
%
    \mathbb{E}_{\text{data}}[\mathbf{x}\vert\mathbf{x}_{t}]+\lambda\big(\mathbb{E}_{\text{data}}[\mathbf{x}\vert\mathbf{x}_{t},\mathbf{c}]-\mathbb{E}_{\text{data}}[\mathbf{x}\vert\mathbf{x}_{t}]\big)
\end{align*} %\mj{what does ``positive" underneath the first term in the parenthesis mean?}
allowing stronger alignment of the sample with the prompt $\mathbf{c}$ via the scale $\lambda$. The purpose of the additional term is to guide the trajectory in the \textit{sharpening direction} toward a desired condition $\mathbf{c}$.

When there are unsafe words in the input text prompt, \textit{SAFREE}~\citep{yoon2024safree} detects unsafe words (tokens) and modifies the unsafe token embeddings. It filters out undesirable concepts with
\begin{align}\label{eq:SAFREE}
    \mathbb{E}_{\text{data}}[\mathbf{x}\vert\mathbf{x}_{t}]+\underbrace{\lambda(\mathbb{E}_{\text{data}}[\mathbf{x}\vert\mathbf{x}_{t},\tilde{\mathbf{c}}_{+}]-\mathbb{E}_{\text{data}}[\mathbf{x}\vert\mathbf{x}_{t}])}_{\text{SAFREE}},
\end{align}
where $\tilde{\mathbf{c}}_{+}$ is a modified prompt embeddings. This altered prompt embedding steers the generation process away from the predefined unsafe concepts.

Another way of negating unsafe concepts is using \textit{negative guidance}~\cite{liu2022compositional}. It reverses the CFG gradient direction for an undesired prompt denoted by $\mathbf{c}_{-}$. 
Formally, one replaces the standard CFG with
\begin{align*}
    \mathbb{E}_{\text{data}}[\mathbf{x}\vert\mathbf{x}_{t}]+\lambda\big(\underbrace{\mathbb{E}_{\text{data}}[\mathbf{x}\vert\mathbf{x}_{t},\mathbf{c}_{+}]}_{\text{positive}}-\underbrace{\mathbb{E}_{\text{data}}[\mathbf{x}\vert\mathbf{x}_{t},\mathbf{c}_{-}]}_{\text{negative}}\big),
\end{align*}
where $\mathbf{c}_{+}$ denotes a positive condition and $\mathbf{c}_{-}$ represents a negative context that we want to avoid.

On the line of negative prompting, \textit{Safe Latent Diffusion} (SLD)~\cite{schramowski2023safe} introduces a guidance by
\begin{align}\label{eq:SLD}
    \mathbb{E}_{\text{data}}[\mathbf{x}\vert\mathbf{x}_{t}]+\underbrace{\lambda(\mathbb{E}_{\text{data}}[\mathbf{x}\vert\mathbf{x}_{t},\mathbf{c}_{+}]-\mathbb{E}_{\text{data}}[\mathbf{x}\vert\mathbf{x}_{t}])}_{\text{CFG}}-\underbrace{\mu(\mathbb{E}_{\text{data}}[\mathbf{x}\vert\mathbf{x}_{t},\tilde{\mathbf{c}}_{-}]-\mathbb{E}_{\text{data}}[\mathbf{x}\vert\mathbf{x}_{t}])}_{\text{SLD}},
\end{align}
where $\tilde{\mathbf{c}}_{-}$ represents a predefined set of unsafe prompts suggested by SLD. Hypothetically, suppose we assume $\mu$ was set as $\lambda$. In that case, the SLD guidance simplifies to a negative guidance $\mathbb{E}_{\text{data}}[\mathbf{x}\vert\mathbf{x}_{t}]+\lambda(\mathbb{E}_{\text{data}}[\mathbf{x}\vert\mathbf{x}_{t},\mathbf{c}_{+}]-\mathbb{E}_{\text{data}}[\mathbf{x}\vert\mathbf{x}_{t},\tilde{\mathbf{c}}_{-}])$. A core difference between SLD and negative guidance is that $\mu$ is adaptive, i.e., $\mu=\mu(\mathbf{c}_{+}, \tilde{\mathbf{c}}_{-}; \gamma, \lambda)$, depending on $\mathbf{c}_{+}$ and $\tilde{\mathbf{c}}_{-}$. This weight is proportional to the norm of the difference between denoisiers: $\Vert \mathbb{E}_{\text{data}}[\mathbf{x} \vert \mathbf{x}_{t}, \mathbf{c}_{+}] - \mathbb{E}_{\text{data}}[\mathbf{x} \vert \mathbf{x}_{t}, \tilde{\mathbf{c}}_{-}] \Vert$. A larger norm suggests that the trajectory is likely to be safe, whereas a smaller norm implies potential unsafety.

% before they are used as a condition of DMs. SAFREE defines an unsafe concept embedding space and then projects the embeddings of unsafe tokens into an orthogonal space of the unsafe embedding space. This altered prompt embedding steers the generation process away from the predefined unsafe concepts.

%Instead of directly using $\mathbf{c}_{US}$ as $\mathbf{c}_{-}$, SLD introduces an adaptive weight $\mu(\mathbf{c}_{+}, \mathbf{c}_{US}; \gamma, \lambda)$ proportional to the denoiser difference norm, defined as $\Vert \mathbb{E}_{\text{data}}[\mathbf{x} \vert \mathbf{x}_{t}, \mathbf{c}_{+}] - \mathbb{E}_{\text{data}}[\mathbf{x} \vert \mathbf{x}_{t}, \mathbf{c}_{US}] \Vert$. The magnitude of this norm serves as an indicator of the proximity of the sampling trajectory to the unsafe region. Specifically, a larger norm suggests that the trajectory is likely to be safe, whereas a smaller norm indicates potential unsafety.

%


%

%
%
%
%



\section{Method}
\label{sec_method}

\begin{figure}[t]%
	\centering
 \begin{subfigure}{0.31\linewidth}
		\centering
		\includegraphics[width=\linewidth]{Figures/Method/moons_10.0_v4.pdf}
		\subcaption{weight $\leftarrow\frac{1}{2}\beta^{*}(\mathbf{x}_{t})$}
        \label{fig:beta_1}
	\end{subfigure}
 \hfil
 \begin{subfigure}{0.31\linewidth}
		\centering
		\includegraphics[width=\linewidth]{Figures/Method/moons_6.5_v4.pdf}
		\subcaption{weight $\leftarrow\beta^{*}(\mathbf{x}_{t})$}
        \label{fig:beta_2}
	\end{subfigure}
 \hfil
 \begin{subfigure}{0.31\linewidth}
		\centering
		\includegraphics[width=\linewidth]{Figures/Method/moons_3.0_v4.pdf}
		\subcaption{weight $\leftarrow2\beta^{*}(\mathbf{x}_{t})$}
        \label{fig:beta_3}
	\end{subfigure}
 \vskip -0.05in
	\caption{Effect of the weight value in \thmref{safe}. (a) If we use half the theoretical weight value, samples generated by our weak safe denoiser also cover the unsafe region (i.e., red dots appearing in the blue area). (b) When we use the theoretical value, the samples avoid unsafe regions while covering the whole safe area. (c) If we penalize more with doubled weight value, the samples not only avoid the unsafe data but also negate the \textit{neighborhood} of unsafe data (i.e., there are no red dots in the black area).}
    \label{fig:beta}
	 \vskip -0.1in
\end{figure}


% The negative prompt $\mathbf{c}_{-}$ or the SLD prompt $\mathbf{c}_{US}$ consist of a limited set of pre-selected words by human users, and may not encompass all images intended to be negated. Consequently, instead of ensuring safety based on text prompts, we introduce a method that guarantees safety based on images, which can be combined with existing text-based safety approaches. Furthermore, while text-based negative guidance can enhance safety, its application lacks a theoretical foundation, thereby offering no guarantees regarding the distribution of the samples.
% We propose constructing a sampling trajectory that adheres to the safe distribution by using a \textit{safe denoiser} defined below to address these issues. 

Text-based prompts (like $\mathbf{c}_{-}$ or $\tilde{\mathbf{c}}_{-}$) rely on limited, user-selected words and may miss undesired content. To address this, we introduce a method that 
offers safety guarantees on generated images, 
which can be combined with existing text-based safety approaches.
% and  complements existing text-based approaches. Below, we start by defining the \textit{safe denoiser}.


\subsection{Safe Denoiser}\label{sec:safe}


We first define an indicator function,  $1_{\text{safe}}(\mathbf{x})$, taking the value of $1$ if $\mathbf{x}$ is safe and $0$ if not. Similarly, we define an indicator function, $1_{\text{unsafe}}(\mathbf{x})$ taking the value of $1$ if $\mathbf{x}$ is unsafe and $0$ if not. 
% Hence, 
These indicator functions are the partition of the unity, resulting in $1= 1_{\text{safe}}(\mathbf{x}) + 1_{\text{unsafe}}(\mathbf{x})$ for all $\mathbf{x}\in \text{supp}(p_{\text{data}})$. Then, we define the following concepts.
% \begin{definition}
%     The unnormalized density of the safe distribution $p_{\text{safe}}(\mathbf{x})$ is $1_{\text{safe}}(\mathbf{x})p_{\text{data}}(\mathbf{x})$. The \textit{safe denoiser} is defined by
%     \begin{align*}
%     \mathbb{E}_{\text{safe}}[\mathbf{x}\vert\mathbf{x}_{t}]=\int \mathbf{x}\frac{p_{\text{safe}}(\mathbf{x})q_{t}(\mathbf{x}_{t}\vert\mathbf{x})}{p_{\text{safe},t}(\mathbf{x}_{t})}\diff\mathbf{x},
% \end{align*}
% where $p_{\text{safe},t}(\mathbf{x}_{t})$ is the marginal distribution of the diffusion process (at time $t$) starting from the safe distribution. Analogously, the unnormalized density of the unsafe distribution $p_{\text{unsafe}}(\mathbf{x})$ is $1_{\text{unsafe}}(\mathbf{x})p_{\text{data}}(\mathbf{x})$. The \textit{unsafe denoiser} is
% \begin{align}\label{eq:E_unsafe}
%     \mathbb{E}_{\text{unsafe}}[\mathbf{x}\vert\mathbf{x}_{t}]=\int \mathbf{x}\frac{p_{\text{unsafe}}(\mathbf{x})q_{t}(\mathbf{x}_{t}\vert\mathbf{x})}{p_{\text{unsafe},t}(\mathbf{x}_{t})}\diff\mathbf{x},
% \end{align}
% where $p_{\text{unsafe},t}(\mathbf{x}_{t})$ is the marginal distribution of the diffusion process (at $t$) starting from the unsafe distribution.
% \end{definition}
\begin{mydefinitionbox}
\noindent\textbf{Definition 3.1.}\label{thm:def} The unnormalized density of the safe distribution $p_{\text{safe}}(\mathbf{x})$ is $1_{\text{safe}}(\mathbf{x})p_{\text{data}}(\mathbf{x})$. The \textit{safe denoiser} is defined by
    \begin{align*}
    \mathbb{E}_{\text{safe}}[\mathbf{x}\vert\mathbf{x}_{t}]=\int \mathbf{x}\frac{p_{\text{safe}}(\mathbf{x})q_{t}(\mathbf{x}_{t}\vert\mathbf{x})}{p_{\text{safe},t}(\mathbf{x}_{t})}\diff\mathbf{x},
\end{align*}
where $p_{\text{safe},t}(\mathbf{x}_{t})$ is the marginal distribution of the diffusion process (at time $t$) starting from the safe distribution. Analogously, the unnormalized density of the unsafe distribution $p_{\text{unsafe}}(\mathbf{x})$ is $1_{\text{unsafe}}(\mathbf{x})p_{\text{data}}(\mathbf{x})$. The \textit{unsafe denoiser} is
\begin{align}\label{eq:E_unsafe}
    \mathbb{E}_{\text{unsafe}}[\mathbf{x}\vert\mathbf{x}_{t}]=\int \mathbf{x}\frac{p_{\text{unsafe}}(\mathbf{x})q_{t}(\mathbf{x}_{t}\vert\mathbf{x})}{p_{\text{unsafe},t}(\mathbf{x}_{t})}\diff\mathbf{x},
\end{align}
where $p_{\text{unsafe},t}(\mathbf{x}_{t})$ is the marginal distribution of the diffusion process (at $t$) starting from the unsafe distribution.
\end{mydefinitionbox}



Our interest is to obtain $\mathbb{E}_{\text{safe}}[\mathbf{x}\vert\mathbf{x}_{t}]$ given the data denoiser $\mathbb{E}_{\text{data}}[\mathbf{x}\vert\mathbf{x}_{t}]$. The theorem below describes the relationship between our safe denoiser and the data denoiser. The proof is given in Appendix \ref{sec:proof}.

\begin{mydefinitionbox}
{% \renewcommand의 범위를 이 그룹으로 제한합니다.
        \renewcommand{\thedummylabelcounter}{3.2}% \ref가 표시할 내용을 "3.2"로 강제 지정
        \refstepcounter{dummylabelcounter}% LaTeX 참조 시스템에 이 "번호"를 알림
        \label{thm:safe}% 레이블 지정 (이제 "3.2"와 연결됨)
    }% 그룹 끝
    \noindent\textbf{Theorem 3.2.} Suppose that $\mathbb{E}_{\textup{data}}[\mathbf{x}\vert\mathbf{x}_{t}]$, $\mathbb{E}_{\textup{safe}}[\mathbf{x}\vert\mathbf{x}_{t}]$, and $\mathbb{E}_{\textup{unsafe}}[\mathbf{x}\vert\mathbf{x}_{t}]$ are the data denoiser, the safe denoiser, and the unsafe denoiser. Then,
    \begin{align}\label{eq:E_safe}
        \mathbb{E}_{\textup{safe}}[\mathbf{x}\vert\mathbf{x}_{t}]=\mathbb{E}_{\textup{data}}[\mathbf{x}\vert\mathbf{x}_{t}]+\beta^{*}(\mathbf{x}_{t})\big(\mathbb{E}_{\textup{data}}[\mathbf{x}\vert\mathbf{x}_{t}]-\mathbb{E}_{\textup{unsafe}}[\mathbf{x}\vert\mathbf{x}_{t}]\big)
    \end{align}
    for a weight defined by
    \begin{align}\label{eq:beta}
        \beta^{*}(\mathbf{x}_{t}) = \frac{Z_{\textup{unsafe}}p_{\textup{unsafe},t}(\mathbf{x}_{t})}{Z_{\textup{safe}}p_{\textup{safe},t}(\mathbf{x}_{t})},
    \end{align}
    where $Z_{\textup{safe}}:=\int 1_{\textup{safe}}(\mathbf{x})p_{\textup{data}}(\mathbf{x})\diff\mathbf{x}$ and $Z_{\textup{unsafe}}:=\int 1_{\textup{unsafe}}(\mathbf{x})p_{\textup{data}}(\mathbf{x})\diff\mathbf{x}$ are normalizing constants of unnormalized densities of safe and unsafe distributions, respectively.
\end{mydefinitionbox}
% \begin{theorem}\label{thm:safe}
% Suppose that $\mathbb{E}_{\textup{data}}[\mathbf{x}\vert\mathbf{x}_{t}]$, $\mathbb{E}_{\textup{safe}}[\mathbf{x}\vert\mathbf{x}_{t}]$, and $\mathbb{E}_{\textup{unsafe}}[\mathbf{x}\vert\mathbf{x}_{t}]$ are the data denoiser, the safe denoiser, and the unsafe denoiser. Then,
%     \begin{align}\label{eq:E_safe}
%         \mathbb{E}_{\textup{safe}}[\mathbf{x}\vert\mathbf{x}_{t}]&=\mathbb{E}_{\textup{data}}[\mathbf{x}\vert\mathbf{x}_{t}]\\
%         &\quad+\beta^{*}(\mathbf{x}_{t})\big(\mathbb{E}_{\textup{data}}[\mathbf{x}\vert\mathbf{x}_{t}]-\mathbb{E}_{\textup{unsafe}}[\mathbf{x}\vert\mathbf{x}_{t}]\big) \nonumber
%     \end{align}
%     for a weight is defined by
%     \begin{align}\label{eq:beta}
%         \beta^{*}(\mathbf{x}_{t}) = \frac{Z_{\textup{unsafe}}p_{\textup{unsafe},t}(\mathbf{x}_{t})}{Z_{\textup{safe}}p_{\textup{safe},t}(\mathbf{x}_{t})},
%     \end{align}
%     where $Z_{\textup{safe}}:=\int 1_{\textup{safe}}(\mathbf{x})p_{\textup{data}}(\mathbf{x})\diff\mathbf{x}$ and $Z_{\textup{unsafe}}:=\int 1_{\textup{unsafe}}(\mathbf{x})p_{\textup{data}}(\mathbf{x})\diff\mathbf{x}$ are normalizing constants of safe and unsafe distributions, respectively.
% \end{theorem}




%
%
%
%
%
%
%
%
%
%
%
%

The theorem above suggests that a safe denoiser can be constructed similarly to CFG. 
%
%
In our case, the denoiser is penalized by $\beta^{*}(\mathbf{x}_{t})$, designed to increase when $\mathbf{x}_{t}$ is likely unsafe. Specifically, a term in the numerator, $p_{\text{unsafe},t}(\mathbf{x}_{t})=\int p_{\text{unsafe}}(\mathbf{x}) q_{t}(\mathbf{x}_{t} \vert \mathbf{x})\diff\mathbf{x}$, grows as the \textit{likelihood of} $\mathbf{x}_{t}$ \textit{being unsafe} increases. In contrast, the denominator grows as the \textit{likelihood of} $\mathbf{x}_{t}$ \textit{being safe} increases. Consequently, $\beta^{*}(\mathbf{x}_{t})$ decreases as $\mathbf{x}_{t}$ becomes more likely to be safe. This indicates that our $\beta^{*}(\mathbf{x}_{t})$ shares a similar intuition to the adaptive weight $\mu$ observed in SLD, but correctly aligns with the intended penalty mechanism. In other words, if $\mathbf{x}_{t}$ is more unsafe than $\mathbf{\tilde{x}}_{t}$, then the trajectory of $\mathbf{x}_{t}$ is more penalized than that of $\mathbf{\tilde{x}}_{t}$, i.e., $\beta^{*}(\mathbf{x}_{t})>\beta^{*}(\mathbf{\tilde{x}}_{t})$.

To provide more intuition on the role of the weight in our theorem, we vary the values that the weight can take and show the corresponding samples. In \figref{beta_1}, we observe that when safety is considered less rigorously than the measure of $\beta{^*}(\mathbf{x}_{t})$, some samples reside within the unsafe region. In contrast, \figref{beta_2} demonstrates that by doubling the safety threshold, both the unsafe region and its immediate surroundings are effectively avoided.
However, in \figref{beta_3}, we observe that the samples from our safe denoiser do not cover the entire safe regions in the data distribution. 

\subsection{Practial Considerations}
\label{sec:unbiased}

For computing \eqref{E_safe}, we need to compute three terms: the data denoiser $\mathbb{E}_{\text{data}}[\mathbf{x}\vert\mathbf{x}_{t}]$, the unsafe denoiser $\mathbb{E}_{\text{unsafe}}[\mathbf{x}\vert\mathbf{x}_{t}]$ and the weight $\beta^{*}(\mathbf{x}_{t})$. We approximate $\mathbb{E}_{\text{data}}[\mathbf{x}\vert\mathbf{x}_{t}]$ by utilizing a pre-trained diffusion model. Consequently, the task reduces to deriving $\mathbb{E}_{\text{unsafe}}[\mathbf{x}\vert\mathbf{x}_{t}]$ and the weight. In this section, we describe our approach to approximating these quantities. 


\paragraph{Approximation of the unsafe denoiser.}
First, we present an approximation of the unsafe denoiser as follows.
%
Given a set of unsafe data points denoted by $\mathbf{x}^{(1)},...,\mathbf{x}^{(N)}$,
\begin{align}\label{eq:est_unsafe}
    \hat{\mathbb{E}}_{\text{unsafe}}[\mathbf{x}\vert\mathbf{x}_{t}]
    %
    =
    %
\sum_{n=1}^{N}\mathbf{x}^{(n)}\frac{q_{t}(\mathbf{x}_{t}\vert\mathbf{x}^{(n)})}{\sum_{m=1}^{N}q_{t}(\mathbf{x}_{t}\vert\mathbf{x}^{(m)})}.
\end{align}
%
%
Each numerator and denominator terms of \eqref{est_unsafe} approximates the numerator and denominator terms of \eqref{E_unsafe}, respectively. It shows that an unsafe denoiser can be expressed as a weighted sum of the unsafe dataset. Here, the weights $\{\frac{q_{t}(\mathbf{x}_{t}\vert\mathbf{x}^{(n)})}{\sum_{m=1}^{N}q_{t}(\mathbf{x}_{t}\vert\mathbf{x}^{(m)})}\}$ form a sum-to-one normalized vector across the unsafe data points, so the unsafe denoiser is approximated as a mixture of unsafe data points.

%

%
%



%

%

%
%
%
%
%
%
%

%


%

%
%
%
%
%
%
%
%
%
%
%
%
%
%
%
%
%

%
%
%
%
%
%



%

%

%

%
%

%

%

%


%
%
%
%
%
%
%
%
%
%
%
%


\paragraph{Approximation of the weight.}
Next, we turn our attention to the computation of $\beta^{*}(\mathbf{x}_{t})$ in \eqref{beta}. Direct calculation is intractable due to the denominator $Z_{\text{safe}}\int p_{\text{safe}}(\mathbf{x})q_{t}(\mathbf{x}_{t}\vert\mathbf{x})$, which is computationally infeasible\footnote{It requies computing $q_{t}(\mathbf{x}_{t}\vert\mathbf{x})$ over all safe data $\mathbf{x}\sim p_{\text{safe}}(\mathbf{x})$, where safe data includes the entire training dataset excluding few unsafe data. Modern text-to-image models like Stable Diffusion~\cite{rombach2022high} are trained with billions of training data~\cite{schuhmann2022laion}, and is infeasible to iterate the entire data at inference time.} to evaluate at every sampling steps. 
%
To address this challenge, 
we approximate $\beta^*$ as
\begin{align*}
    \beta^*(\mathbf{x}_{t}) \approx \eta \cdot \beta(\mathbf{x}_{t}),
\end{align*}
with a constant $\eta$ and a function $\beta(\mathbf{x}_{t})$ defined by
%
\begin{align*}
    \beta(\mathbf{x}_{t})&=\int p_{\text{unsafe}}(\mathbf{x})q_{t}(\mathbf{x}_{t}\vert\mathbf{x})\diff\mathbf{x}\approx \frac{1}{N}\sum_{n=1}^{N}q_{t}(\mathbf{x}_{t}\vert\mathbf{x}^{(n)}) 
\end{align*}
where the last line is an unbiased estimate of $\beta$. We treat $\eta$ as a controllable hyperparmeter, with which we replace the computation of the remaining terms in \eqref{beta}. 
%
This approximation is reasonable insofar as the numerator alone captures the overall trend of $\beta^{*}(\mathbf{x}_{t})$: as $\mathbf{x}_{t}$ becomes more likely to be unsafe, both $\beta^{*}(\mathbf{x}_{t})$ and the numerator increase correspondingly. This approximation of the weight significantly reduces computational complexity. Additionally, we observe that applying the safe denoiser at the final stage of sampling (i.e., when $t$ is small) hurts the sample quality, since the signal from unsafe denoiser--a weighted sum of unsafe data points--acts as a structural noise for detailed denoising. From this observation, we propose to apply the safe denoiser only at the beginning of sampling process.
%

 %
%
%
%
%

\input{Manuscript/sources/algorithms}

\paragraph{Putting things together.}
With these approximations mentioned above, we arrive at the final safe denoiser:
\begin{align}\label{eq:uncond_safe}
\begin{split}
   \hat{\mathbb{E}}_{\text{safe}}[\mathbf{x}\vert\mathbf{x}_{t}] = \mathbb{E}_{\text{data}}[\mathbf{x}\vert\mathbf{x}_{t}]
+ \eta\beta(\mathbf{x}_{t})(\mathbb{E}_{\text{data}}[\mathbf{x}\vert\mathbf{x}_{t}]-\hat{\mathbb{E}}_{\text{unsafe}}[\mathbf{x}\vert\mathbf{x}_{t}]),
\end{split}
\end{align}
where $\hat{\mathbb{E}}$ is given in \eqref{est_unsafe}. 
%
Our results in \secref{experiments} validate the effectiveness of our approximations in ensuring sample safety without incurring prohibitive computational costs.
%

%

%
%
%
%
%
%



%
%
%
%
%
%

%

%

%

\subsection{Extending Safe Denoiser to Text-to-Image generation}

\begin{wraptable}{r}{0.32\columnwidth}
\vskip -0.3in
\caption{Joint effect of existing text-based guidance (SAFREE) and ours. We evaluate the attack success rate. Both "No" with 0.962 refers to SD-v1.4~\cite{rombach2022high} with CFG. The lower, the better.}
% \vskip -0.1in
  \label{tab:sole}
  \centering
\begin{tabular}{lccc}
% \xmark
% \cmark
\toprule
&& \multicolumn{2}{c}{Neg. Prompt} \\
&& No & Yes \\\midrule
\multirow{2}{*}{Ours} & No & 0.962 & 0.601 \\
& \cc{15}Yes & \cc{15}0.633 & \cc{15}\textbf{0.469} \\
\bottomrule
\end{tabular}
\end{wraptable}
%
% While our methodology is effective as a standalone algorithm, 
% % it operates orthogonally to existing text-based safety mechanisms within text-conditional generation frameworks. This inherent orthogonality facilitates its straightforward 
% we can straightforwardly
% integrate our method, as a plug-in component, into established text-based safety mechanisms, enhancing the safety level as presented in 
While our methodology is effective as a standalone algorithm, we can also integrate it straightforwardly as a plug-in component into established text-based safety mechanisms, thereby enhancing the overall safety level, as shown in \tabref{sole} \footnote{We tested on MMA-Diffusion~\citep{yang2024mma} nudity prompts and measure the rate the model generates unsafe images. See Section~\ref{sec:experiments} and Table~\ref{tab:t2i} for further details.}. For example, when our approach is combined with SAFREE, the predicted clean sample $\mathbf{x}_{0\vert t}$ (representing the estimated data at step $t=0$ given a sample $\mathbf{x}_{t}$ at step $t$) can be computed by
\begin{align}\label{eq:SAFREE_ours}
\mathbf{x}_{0\vert t}=\mathbb{E}_{\text{safe}}[\mathbf{x}\vert\mathbf{x}_{t}]
+\underbrace{\lambda(\mathbb{E}_{\text{data}}[\mathbf{x}\vert\mathbf{x}_{t},\tilde{\mathbf{c}}_{+}]-\mathbb{E}_{\text{data}}[\mathbf{x}\vert\mathbf{x}_{t}])}_{\text{SAFREE}}.
\end{align}
When it is combined with SLD, the formula is as follows:
\begin{align}\label{eq:SLD_ours}
\mathbf{x}_{0\vert t}=\mathbb{E}_{\text{safe}}[\mathbf{x}\vert\mathbf{x}_{t}]
+\underbrace{\lambda(\mathbb{E}_{\text{data}}[\mathbf{x}\vert\mathbf{x}_{t},\mathbf{c}_{+}]-\mathbb{E}_{\text{data}}[\mathbf{x}\vert\mathbf{x}_{t}])}_{\text{CFG}}-\underbrace{\mu(\mathbb{E}_{\text{data}}[\mathbf{x}\vert\mathbf{x}_{t},\tilde{\mathbf{c}}_{-}]-\mathbb{E}_{\text{data}}[\mathbf{x}\vert\mathbf{x}_{t}])}_{\text{SLD}}.
\end{align}
Note these \eqref{SAFREE_ours} and \eqref{SLD_ours} replaces the data denoiser $\mathbb{E}_{\text{data}}[\mathbf{x}\vert\mathbf{x}_{t}]$ by the safe denoiser $\mathbb{E}_{\text{safe}}[\mathbf{x}\vert\mathbf{x}_{t}]$, compared to \eqref{SAFREE} and \eqref{SLD}, respectively. In implementation, as described in \secref{unbiased}, we approximate the safe denoiser by \eqref{uncond_safe}. In diffusion sampling, we utilize this safe $\mathbf{x}_{0\vert t}$ in either DDPM~\cite{ho2020denoising} or DDIM~\cite{song2020denoising}, see \algoref{safer} for details.

% Using this denoiser allows us to negate data samples based on the information from the images (from our safe denoiser) and the information based on the prompts (from both CFG and SLD). Note this \eqref{safer} includes only the additional term for the safe denoiser compared to \eqref{SLD}. In implementation, as described in \secref{unbiased}, we approximate the second term of \eqref{safer} by \eqref{uncond_safe}. In diffusion sampling, we utilize this safe $\mathbf{x}_{0\vert t}$ in either DDPM~\cite{ho2020denoising} or DDIM~\cite{song2020denoising}, see \algoref{safer} for details.

When our safe denoiser is combined  with the text-based guidance methods, we introduce a new set of hyperparameters $\beta_{t}$, such that we set $\beta(\mathbf{x}_{t})$ to zero if this value falls below a predefined threshold $\beta_{t}$. This condition indicates that if a sample $\mathbf{x}_{t}$ is sufficiently safe, modifying the trajectory is no longer necessary. This thresholding improves accuracy thanks to their better controllability relative to the text guidance terms.

%
%
%
%





%

%
%
%
%

\section{Related Work}\label{sec:related}
%
%
%
%
%

Earlier work on machine unlearning in generative modelling focused on object unlearning in classification (forgetting images from a selected class), unconditional image generation (forgetting harmful images) or concept erasing (forgetting harmful concepts). Most of the work belonging to this category required retraining the entire generative models or some part of them, rather than modifying the sampling trajectory or input prompts \cite{heng2023selective, li2024machine, adapt_then_unlearn, zhang2024forgetmenot, gandikota2023unified, lu2024mace, gong2024reliable, lu2024mace}. 
In more recent work, training-free and text-based methods have also emerged as computationally efficient alternatives \cite{schramowski2023safe,yoon2024safree,ban2024understanding, armandpour2023reimagine}. 
However, most of these approaches lack a theoretical ground, unlike our work. 


%
%

%

%

%
%

%

%

%
%
%
%

 





%


Despite these advances, generative models remain susceptible to adversarial prompts, malicious manipulations of learnable parameters, textual cues, or even random noise \cite{pham2023circumventing, chin2024promptingdebugging, zhang2024generate, tsai2024ringabell}. These findings highlight using a single defense such as concept erasing as a standalone solution may be insufficient to ensure safe content generation. We see this as an opportunity for our method to be combined with powerful text-based defense mechanisms to enhance their performance.  



%

A closely related recent work, \textit{Sparse Repellency} (SR)~\cite{kirchhof2024sparse}, is a training-free technique that modifies the denoising trajectory to avoid unsafe images. Their denoiser follows $\mathbb{E}_{\text{data}}[\mathbf{x}\vert\mathbf{x}_{t}]+\sum_{n=1}^{N}\text{ReLU}\left(\tfrac{r}{\Vert\mathbb{E}_{\text{data}}[\mathbf{x}\vert\mathbf{x}_{t}]-\mathbf{x}^{(n)}\Vert}-1\right) \times(\mathbb{E}_{\text{data}}[\mathbf{x}\vert\mathbf{x}_{t}]-\mathbf{x}^{(n)}).$
%
%
%
ReLU activation ensures that the diffusion trajectory is penalized when the denoiser falls within the neighborhood of radius $r$ around unsafe data, and remains unmodified otherwise. 
Given a single unsafe image, $\text{ReLU}\left(\tfrac{r}{\Vert\mathbb{E}_{\text{data}}[\mathbf{x}\vert\mathbf{x}_{t}]-\mathbf{x}^{(n)}\Vert}-1\right) (\mathbb{E}_{\text{data}}[\mathbf{x}\vert \mathbf{x}_{t}]-\mathbf{x}^{(n)})$ resembles the second term, $\mathbb{E}_{\text{safe}}[\mathbf{x}\vert\mathbf{x}_{t}]$, in \eqref{E_safe} if the ReLU activation is comparable to our $\beta^*$. From this point of view, our method can be regarded as a generalization of SR.
However, unlike our method, SR does not guarantee that the samples are from a safe distribution. 

% \textit{Diffusion Soup}~\cite{biggs2024diffusion} shares a similar theoretical analysis to ours. They propose merging the weights of DMs separately trained on independent subsets of data, resulting in a sampling distribution that extends beyond the framework outlined in our \thmref{safe}. Unlike their method, we do not require additional fine-tuning, and our analysis defines the safe and unsafe denoisers and their explicit relationship between those. Recently, \textit{Dynamic Negative Guidance} (DNG)~\cite{koulischer2025dynamic} also addresses the strength of negative guidance using a similar theoretical framework. While DNG relies on sequential computation based on a Markov chain, our approach estimates denoisers for safe and unsafe in an expectation manner, which avoids sequential dependencies.
% Additionally, DNG requires extra training for unsafe denoiser in cases where text-to-image generation is not applicable, whereas our method is entirely training-free. 
\textit{Diffusion Soup}~\cite{biggs2024diffusion} presents a related theoretical analysis by merging DMs trained on separate data subsets, but requires fine-tuning. In contrast, our method is training-free and formally defines safe and unsafe denoisers and their relationship. \textit{Dynamic Negative Guidance} (DNG)~\cite{koulischer2025dynamic} also uses a similar framework but relies on sequential computation based on a Markov chain and requires extra training for unsafe denoisers, whereas our approach estimates safe and unsafe denoisers in an expectation manner without training overhead.
\section{Experiments}
\label{sec:experiments}


We present the experimental results of our method, Safe Denoiser. Section~\ref{sec:t2i} details the outcomes of our text-to-image generation experiments, while the subsequent section explores both class-conditional and unconditional image generation.

\subsection{Text-to-Image Generation}
\label{sec:t2i}

In this section, we conduct an in-depth analysis of safety issues in text-to-image models, focusing on tasks involving nudity and inappropriate content. The nudity task evaluates how well safety methods prevent harmful outputs as attack difficulty gradually increases. In contrast, the inappropriate content task examines whether these safety methods remain effective when handling multiple concepts simultaneously. 
We use \textit{Stable Diffusion} (SD)~\cite{rombach2022high} v1.4\footnote{\url{https://huggingface.co/CompVis/stable-diffusion-v1-4}} with DDPM sampler. 
To evaluate safety, we follow previous studies by assessing Attack Success Rate (ASR), Toxic Rate (TR), and Inappropriate Probability (IP)~\cite{schramowski2023safe, yoon2024safree}. 
We measure ASR by the proportion of generated images that exceeds 0.6 nude class probability, measured by NudeNet\footnote{\url{https://github.com/notAI-tech/NudeNet}}. The TR is computed by the average of nude class probability, measured also by NudeNet. The IP is the classification probability score of generating inappropriate images, measured by the Q16 classifier~\citep{schramowski2022can}. For the nudity task, we select 515 unsafe images from I2P~\cite{schramowski2023safe} that exceeds 0.6 nude class probability. For the inappropriate content tasks, we randomly sample 3,000 images from I2P as the unsafe dataset. To evaluate, we use the broder dataset CoPro~\cite{liu2024latent}, which covers the same categories of I2P. Notably, all experiments uses the identical unsafe datasets across all baselines for consistency, see Appendix~\ref{sec:t2i_detail} for details. 

Besides safety-related metrics, we prioritize maintaining high image quality and prompt alignment simultaneously. To this end, we calculate Fr\'echet Inception Distance (FID)~\cite{heusel2017gans} for the generation fidelity and CLIP~\cite{radford2021learning} to measure whether the samples follow human instructions. We use a PyTorch package~\cite{Seitzer2020FID} to compute the FID by comparing 10K reference images selected from the COCO-2014~\cite{lin2014microsoft} validation split and 10K generated images from the prompts identically selected from the same COCO dataset. Also, we evaluate the CLIP score~\citep{radford2021learning} using ViT-B-32\footnote{\url{https://huggingface.co/openai/clip-vit-base-patch32}}.

\paragraph{Safe Generation against Nudity Prompts} 
\tabref{t2i} summarizes our experimental findings. In these experiments, we utilize unsafe prompts proposed by Ring-A-Bell~\cite{tsai2024ringabell} (79 prompts), UnlearnDiff~\cite{zhang2024generate} (116 sexual prompts), and MMA-Diffusion~\cite{yang2024mma} (1000 prompts). These prompts are adversarially generated to fool the existing generative models. For baseline comparisons, we consider both training-based approaches, specifically \textit{ESD}~\cite{gandikota2023erasing} and \textit{RECE}~\cite{gong2024reliable}, and training-free methods such as SLD~\cite{schramowski2023safe} and SAFREE~\cite{yoon2024safree}. Initially, we observe that about $96.2\%$ of generated SD-v1.4 images are unsafe when using MMA-Diffusion prompts. Existing baselines demonstrate performance improvements over SD across datasets. 



% \begin{figure*}[!t]
%     \centering
%     \includegraphics[width=0.99\textwidth]{Figures/Experiments/imagenet/vs_dogs/comparison_imagenet_chihuahua_vs_dogs.pdf}
%     \caption{Generated samples when negating the Chihuahua class, primarily producing visually similar small dog breeds.}
%     \label{fig:imagenet_comparison_dogs}
% \end{figure*}

% \begin{table*}[!t]
%
% \begin{minipage}[t]{.45\textwidth}
% \caption{Performance evaluation in FFHQ. We use ResNet18~\cite{he2016deep} to classify the sex of generated samples. We compute FID by comparing 1K male subset of CelebA validation and 1K generated images.}
%     \label{tab:ffhq}
%     \centering
%     \resizebox{0.95\textwidth}{!}{%
%     \begin{tabular}{lccc}
%         \toprule
%         Method & Female & Male & FID $\downarrow$ \\\midrule
%         Baseline (B)         & $64.0\%$ & $36.0\%$ & 109.07  \\
%         B + SR   & $53.1\%$ & $46.9\%$ & 130.52   \\
%         \cc{15}B + Ours       & \cc{15}$55.6\%$ & \cc{15}$44.4\%$ & \cc{15}96.57   \\
%         \bottomrule
%     \end{tabular}
%     }%
%     \end{minipage}%
%     \hfill
%     \begin{minipage}[t]{.53\textwidth}
%     \caption{Experiments on ImageNet for the specific class (Chihuahua) negation task. Top-1 is the classification accuracy of the generated samples on 999 classes, and Top-1* indicates the accuarcy on the specific class.} \label{tab:imagenet} 
%     \centering 
%     \resizebox{0.99\textwidth}{!}{%
%     \begin{tabular}{lccccc} 
%         \toprule 
%         Method & Prec $\uparrow$ & Rec $\uparrow$ & Top-1 $\uparrow$ & Top-1$^{*}$ $\downarrow$ \\ 
%         \midrule 
%         Baseline (B) & 0.72 & 0.63 & 0.76 & 0.68 \\ 
%         B + SR & 0.59 & 0.54 & 0.01 & 0.0 \\ 
%         \cc{15}B + Ours & \cc{15}0.62 & \cc{15}0.58 & \cc{15}0.14 & \cc{15}0.0 \\ \bottomrule 
%     \end{tabular} 
%     }%
%     \end{minipage}%
% \end{table*}

Our method, combined with SLD or SAFREE, significantly improves safety performance while maintaining image quality. Notably, the extent of improvement varies considerably depending on the characteristics of the prompts. For instance, with MMA-Diffusion prompts, %
the performance of text-based baselines (like SLD) is markedly inferior ($88.1\%$ generated images are unsafe) compared to their performance on other prompt datasets such as Ring-A-Bell or UnlearnDiff. This discrepancy arises because MMA-Diffusion prompts lack explicit nudity information due to being part of a white-box adversarial attack, making it challenging for text-based safety methods to erase such concepts. In contrast, our approach employs purely image-based guidance, which results in substantial performance gains from $88.1\%$ to $48.1\%$ in ASR on MMA-Diffusion when combined with existing text-based methods. Our method significantly improves the performance across all other prompt datasets, not limited to MMA-Diffusion.

% New Contents in NeurIPS2025

\begin{table}[!t]
    \caption{Performance comparison of baselines on various datasets in safe generation against nudity prompts. Our method, combined with existing approaches, significantly improves the safety performance while keeping image quality.}
    \label{tab:t2i}
    \centering
    \resizebox{0.99\textwidth}{!}{%
        \begin{tabular}{lccccccccccc}
            \toprule
            \multirow{2}{*}{Method} 
            & \multirow{2}{*}{\shortstack{Fine\\Tuning}} & \multirow{2}{*}{\shortstack{Negative\\Prompt}} & \multirow{2}{*}{\shortstack{Safe\\Denoiser}}
            & \multicolumn{2}{c}{Ring-A-Bell} 
            & \multicolumn{2}{c}{UnlearnDiff}  
            & \multicolumn{2}{c}{MMA-Diffusion}  
            & \multicolumn{2}{c}{COCO-30K} \\
            \cmidrule(lr){5-6} \cmidrule(lr){7-8} \cmidrule(lr){9-10} \cmidrule(lr){11-12}
            & 
            & 
            & 
            & ASR $\downarrow$ & TR $\downarrow$ 
            & ASR $\downarrow$ & TR $\downarrow$ 
            & ASR $\downarrow$ & TR $\downarrow$ 
            & FID $\downarrow$ & CLIP $\uparrow$ \\\midrule
            SD-v1.4 & - & - & - & 0.797 & 0.809 & 0.809 & 0.845 & 0.962 & 0.956 & 25.04 & \textbf{31.38} \\\midrule
            ESD & \cmark & \xmark & \xmark & 0.456 & 0.506 & 0.422 & 0.426 & 0.628 & 0.640 & 27.38 & 30.59 \\
            RECE & \cmark & \xmark & \xmark & 0.177 & 0.212 & 0.284 & 0.292 & 0.651 & 0.664 & 33.94 & 30.29\\%\hline
            % SLD-Weak& \cmark & \cmark & \xmark & 0.785 & 0.774 & 0.828 & 0.791 & 0.953 & 0.949 & 26.60 & 31.30 \\
            % SLD-Medium & \cmark & \cmark & \xmark & 0.646 & 0.675 & 0.750 & 0.717 & 0.939 & 0.929 & 29.10 & 30.76 \\
            % SLD-Strong & \cmark & \cmark & \xmark & 0.620 & 0.636 & 0.664 & 0.651 & 0.905 & 0.908 & 31.28 & 30.00 \\
            SLD & \xmark & \cmark & \xmark & 0.481 & 0.573 & 0.629 & 0.586 & 0.881 & 0.882 & 36.47 & 29.28 \\
            \cc{15}$$ + Ours & \cc{15}\xmark & \cc{15}\cmark & \cc{15}\cmark & \cc{15}0.354 & \cc{15}0.429 & \cc{15}0.526 & \cc{15}0.485 & \cc{15}0.481 & \cc{15}0.549 & \cc{15}36.59 & \cc{15}29.10 \\
            SAFREE & \xmark & \cmark & \xmark & 0.278 & 0.311 & 0.353 & 0.363 & 0.601 & 0.618 & 25.29 & 30.98 \\%\hline
            \cc{15}$$ + Ours & \cc{15}\xmark & \cc{15}\cmark & \cc{15}\cmark & \cc{15}\textbf{0.127} & \cc{15}\textbf{0.169} & \cc{15}\textbf{0.207} & \cc{15}\textbf{0.241} & \cc{15}\textbf{0.469} & \cc{15}\textbf{0.501} & \cc{15}\textbf{22.55} & \cc{15}30.66 \\  
            % (+ \textit{Safe Denoiser}) &&&&&&&&&&&\\
            % SLD-Weak & \cmark & \cmark & \cmark & 0.557 & 0.616 & 0.690 & 0.669 & 0.805 & 0.807 & 26.23 & 30.50 \\
            % % SLD-Medium & \cmark & \cmark & \cmark & 0.595 & 0.603 & 0.690 & 0.654 & 0.849 & 0.851 & - & - \\
            % SLD-Medium & \cmark & \cmark & \cmark & 0.405 & 0.482 & 0.638 & 0.591 & 0.735 & 0.748 & - & \textcolor{gray}{30.90} \\
            % SLD-Strong & \cmark & \cmark & \cmark & 0.430 & 0.513 & 0.629 & 0.581 & 0.642 & 0.660 & - & \textcolor{gray}{29.50} \\
            % SLD + Ours & \cmark & \cmark & \cmark & 0.354 & 0.429 & 0.526 & 0.485 & 0.481 & 0.549 & - & \textcolor{gray}{28.00} \\
            % % SAFREE & \cmark & \cmark & \cmark & 0.139 & 0.198 & 0.241 & 0.306 & 0.557 & 0.591 & 25.45 & 30.06\\
            % SAFREE + Ours & \cmark & \cmark & \cmark & 0.127 & 0.169 & 0.207 & 0.241 & 0.469 & 0.501 & - & - \\
            \bottomrule
        \end{tabular}
    }%
\end{table}

\begin{table}[!t]
    \caption{Performance of inappropriate probability (IP) and CLIP Score on the CoPro dataset. Our method incoporating with negative prompts enhances safety performance even across multiple concepts simultaneously.}
    \label{tab:ip_CoPro}
    \centering
    \resizebox{0.99\textwidth}{!}{%
    \begin{tabular}{lccccccccc}
        \toprule
        \multicolumn{1}{c}{Method} & \begin{tabular}[c]{@{}c@{}}Harra- \\ sment $\downarrow$\end{tabular} & Hate $\downarrow$  & \begin{tabular}[c]{@{}c@{}}Illegal \\ Activity $\downarrow$\end{tabular} & \begin{tabular}[c]{@{}c@{}}Self- \\ harm $\downarrow$\end{tabular}& Sexual $\downarrow$ & \begin{tabular}[c]{@{}c@{}}Shock- \\ ing $\downarrow$\end{tabular}& \begin{tabular}[c]{@{}c@{}}Viole- \\ nce $\downarrow$\end{tabular}& \begin{tabular}[c]{@{}c@{}}Avg. \\ IP $\downarrow$\end{tabular}  & \multicolumn{1}{c}{CLIP $\uparrow$}  \\
        \cmidrule(lr){1-9} \cmidrule(lr){10-10} 
        SD-v1.4                         & 0.269      & 0.154 & 0.206                                                       & 0.319     & 0.120  & 0.221    & 0.274    & 0.223 & 29.81 \\
        \cc{15} + Ours                    & \cc{15}0.206      & \cc{15}0.148 & \cc{15}0.197                                                       & \cc{15}0.209     & \cc{15}0.109  & \cc{15}0.209    & \cc{15}0.230    & \cc{15}0.187 & \cc{15}29.21     \\\cmidrule(lr){1-10}
        SLD                & 0.223      & 0.106 & 0.161                                                       & 0.247     & 0.078  & 0.158    & 0.217    & 0.170 &29.65  \\
        \cc{15} + Ours                & \cc{15}0.168      & \cc{15}0.113 & \cc{15}0.152                                                       & \cc{15}0.169     & \cc{15}0.078  & \cc{15}0.165    & \cc{15}0.212    & \cc{15}0.151 & \cc{15}28.95 \\\cmidrule(lr){1-10}
        SAFREE                     & 0.182      & 0.118 & 0.144                                                       & 0.183     & 0.085  & 0.150    & 0.206    & 0.153 & 28.91 \\
        \cc{15} + Ours                & \cc{15}0.156      & \cc{15}0.112 & \cc{15}0.161                                                       & \cc{15}0.153     & \cc{15}0.083  & \cc{15}0.159    & \cc{15}0.185    & \cc{15}0.144 & \cc{15}28.49 \\
        \bottomrule
        \end{tabular}
    }%
\end{table}

% \begin{table*}[!t]
%     \caption{Attack Success Rate (ASR) in Nudity}
%     \label{tab:t2i}
%     \centering
%     \resizebox{0.95\textwidth}{!}{%
%         \begin{tabular}{lccc|cccccc|cc}
%             \toprule
%             \multirow{2}{*}{\centering Method} 
%             & \multicolumn{3}{c|}{Modification} 
%             & \multicolumn{2}{c}{Ring-A-Bell} 
%             & \multicolumn{2}{c}{UnlearnDiff}  
%             & \multicolumn{2}{c}{MMA-Diffusion}  
%             & \multicolumn{2}{c}{COCO-30K} \\
%             \cmidrule(lr){2-4} \cmidrule(lr){5-6} \cmidrule(lr){7-8} \cmidrule(lr){9-10} \cmidrule(lr){11-12}
%             & \multirow{2}{*}{\shortstack{Training \\ Free}} 
%             & \multirow{2}{*}{\shortstack{Negative \\ Prompt}} 
%             & \multirow{2}{*}{\shortstack{Image \\ Repellency}} 
%             & ASR $\downarrow$ & Toxic Rate $\downarrow$ 
%             & ASR $\downarrow$ & Toxic Rate $\downarrow$ 
%             & ASR $\downarrow$ & Toxic Rate $\downarrow$ 
%             & FID $\downarrow$ & CLIP $\uparrow$ \\ 
%             &  &  &  &  &  &  &  &  &  &  &  \\ 
%             \midrule
%             SD-v1.4                               & -                                           & -                                      & -                                    & 0.797                                             & 0.809                                                         & 0.809                                                 & 0.845                                                         & 0.962                                                     & 0.956                                                         & 25.04                                     & 31.38                                 \\
%             \hline
%             ESD                                   & \xmark                                      & \xmark                                 & \xmark                               & 0.456                                             & 0.506                                                         & 0.422                                                 & 0.426                                                         & 0.628                                                     & 0.640                                                         & 27.38                                     & 30.59                                 \\
%             RECE                                  & \xmark                                      & \cmark                                 & \xmark                               & 0.177                                             & 0.212                                                         & 0.284                                                 & 0.292                                                         & 0.651                                                     & 0.664                                                         & 33.94                                     & 30.29                                 \\
%             \hline
%             SLD-Weak                              & \cmark                                      & \cmark                                 & \xmark                               & 0.785                                             & 0.774                                                         & 0.828                                                 & 0.791                                                         & 0.953                                                     & 0.949                                                         & 26.60                                     & 31.30                                 \\
%             SLD-Medium                            & \cmark                                      & \cmark                                 & \xmark                               & 0.646                                             & 0.675                                                         & 0.750                                                 & 0.717                                                         & 0.939                                                     & 0.929                                                         & 29.10                                     & 30.76                                 \\
%             SLD-Strong                            & \cmark                                      & \cmark                                 & \xmark                               & 0.620                                             & 0.636                                                         & 0.664                                                 & 0.651                                                         & 0.905                                                     & 0.908                                                         & 31.28                                     & 30.00                                 \\
%             SLD-Max                               & \cmark                                      & \cmark                                 & \xmark                               & 0.481                                             & 0.573                                                         & 0.629                                                 & 0.586                                                         & 0.881                                                     & 0.882                                                         & -                                         & \textcolor{gray}{29.40}                                     \\
%             SAFREE                                & \cmark                                      & \cmark                                 & \xmark                               & 0.278                                             & 0.311                                                         & 0.353                                                 & 0.363                                                         & 0.601                                                     & 0.618                                                         & 25.29                                     & 30.98                                 \\
%             \hline
%             (\textit{Ours;$\beta = 0.03$})        &                                             &                                        &                                      &                                                   &                                                               &                                                       &                                                               &                                                           &                                                               &                                           &                                       \\
%             SLD-Weak                              & \cmark                                      & \cmark                                 & \cmark                               & 0.557                                             & 0.616                                                         & 0.690                                                 & 0.669                                                         & 0.805                                                     & 0.807                                                         & 26.23                                     & 30.50                                 \\
%             SLD-Medium                            & \cmark                                      & \cmark                                 & \cmark                               & 0.405                                             & 0.482                                                         & 0.638                                                 & 0.591                                                         & 0.735                                                     & 0.748                                                         & -                                         & \textcolor{gray}{30.90}                                     \\
%             SLD-Strong                            & \cmark                                      & \cmark                                 & \cmark                               & 0.430                                             & 0.513                                                         & 0.629                                                 & 0.581                                                         & 0.642                                                     & 0.660                                                         & -                                         & \textcolor{gray}{29.50}                                     \\
%             SLD-Max                               & \cmark                                      & \cmark                                 & \cmark                               & 0.354                                             & 0.429                                                         & 0.526                                                 & 0.485                                                         & 0.481                                                     & 0.549                                                         & -                                         & \textcolor{gray}{28.00}                                     \\
%             SAFREE                                & \cmark                                      & \cmark                                 & \cmark                               & 0.139                                             & 0.198                                                         & 0.241                                                 & 0.306                                                         & 0.557                                                     & 0.591                                                         & 25.45                                     & 30.06                                 \\
%             \bottomrule
%         \end{tabular}
%     }%
%     \vspace{-0.07in}
% \end{table*}

% New Contents in NeurIPS2025
\paragraph{Inappropriate Probability in CoPro Dataset}

\tabref{ip_CoPro} presents that our method consistently achieves enhancement of safe content generation against multiple categories while maintaining a balance in textual prompt alignments across all baselines. Since training-based approaches do not provide official checkpoints for this task, we focus on training-free approaches. Overall, our method effectively improves inappropriate probability (IP) on CoPro dataset. Notably, all baselines show a reduction in average IP when combined with our method. Additionally, our method effectively preserves the alignment between human instructions (prompts) and generated images, with any introduced misalignment being minimal, as demonstrated by CLIP scores. Furthurmore, our method performs on-par with previous methods in terms of the sample-wise aesthetic scores, showing that there is a minimal impact in the sample quality by applying our method, see Appendix~\ref{sec:additional_result}. %Our method accomplishes a balanced reduction in IP, avoding unfavorable situations where improvement is limited to specific categories and the others do not benefit. 
These results suggest that our method effectively manages multiple concepts simultaneously while reliably generations away from unsafe content.

\paragraph{Ablation Studies} 

We present a pair of ablation studies to evaluate the robustness and effectiveness of our method. First, \figref{ablation_1} shows the effect of the number of unsafe data points on model performance. We observe that increasing the number of unsafe data points leads to better performance. %

% \begin{figure}[t]%
% 	\centering
%  \begin{subfigure}{0.48\linewidth}
% 		\centering
% 		\includegraphics[width=\linewidth]{Figures/Experiments/unsafe_data.pdf}
%         \vskip -0.05in
% 		\subcaption{Effect on $N$}
% 	\end{subfigure}
%     \hfil
%  \begin{subfigure}{0.48\linewidth}
% 		\centering
% 		\includegraphics[width=\linewidth]{Figures/Experiments/beta.pdf}
%         \vskip -0.05in
% 		\subcaption{Effect on $\beta_{t}$}
% 	\end{subfigure}
%     \hfil
%  % \begin{subfigure}{0.31\linewidth}
% 	% 	\centering
% 	% 	\includegraphics[width=\linewidth]{Figures/Experiments/sld.pdf}
%  %        \vskip -0.05in
% 	% 	\subcaption{Effect on prompt weights}
% 	% \end{subfigure}
%  \vskip -0.05in
% 	\caption{Ablation studies of (a) the effect on the number of unsafe data ($N$), (b) the effect on the threshold ($\beta_{t}$).} %, and (c) the effect on the prompt weights. All metrics are evalauted on MMA-Diffusion}
%     \label{fig:ablation}
% 	 %
% \end{figure}

\begin{figure}[t]
    \begin{minipage}[t]{0.64\linewidth}
        \centering
        \begin{subfigure}{0.48\linewidth}
            \centering
            \includegraphics[width=\linewidth]{Figures/Experiments/unsafe_data.pdf}
            \vskip -0.05in
            \subcaption{Effect on $N$}
            \label{fig:ablation_1}
        \end{subfigure}
        \begin{subfigure}{0.48\linewidth}
            \centering
            \includegraphics[width=\linewidth]{Figures/Experiments/beta.pdf}
            \vskip -0.05in
            \subcaption{Effect on $\beta_{t}$}
            \label{fig:ablation_2}
        \end{subfigure}
        \vskip -0.05in
        \captionof{figure}{Ablation studies of (a) the effect on the number of unsafe data ($N$), (b) the effect on the threshold ($\beta_{t}$).}
        \label{fig:ablation}
    \end{minipage}
    \hfill
    \centering
    \begin{minipage}[t]{0.32\linewidth}
    \vspace*{-3.2cm}
        \centering
        \captionof{table}{Experiments of the class negation on ImageNet. Top-1* is the classification accuracy of the generated samples on the negated class (Chihuahua). Refer Appendix~\ref{sec:pixel_detail} and \ref{sec:additional_result}.}
        \label{tab:imagenet_side_by_side}
        \resizebox{\linewidth}{!}{
            \begin{tabular}{lcccc}
                \toprule
                Method & Prec $\uparrow$ & Rec $\uparrow$ & Top-1* $\downarrow$ \\
                \midrule
                Baseline & 0.72 & 0.63 & 0.68 \\
                B + SR & 0.59 & 0.54 & 0.00 \\
                \cc{15}B + Ours & \cc{15}0.62 & \cc{15}0.58 & \cc{15}0.00 \\ \bottomrule
            \end{tabular}
        }
    \end{minipage}
    \vskip -0.1in
\end{figure}


We then explore the influence of the threshold parameter $\beta_{t}$ (see \algoref{safer}), which governs the application of the safe denoiser. For simplicity, we fixed $\beta_{t}$ across all time steps. \figref{ablation_2} shows the performance exhibits a U-shaped relationship to $\beta_{t}$. Specifically, when $\beta_{t} = 0$, the safe denoiser is applied to all samples $\mathbf{x}_{t}$ regardless of their safety status. Conversely, when $\beta_{t} = \infty$, the safe denoiser is not applied. At intermediate values of $\beta_{t}$, the safe denoiser is applied selectively to a certain proportion of unsafe samples $\mathbf{x}_{t}$. The U-shaped trend indicates 
%
that selectively applying the safe denoiser to unsafe samples based on an appropriate $\beta_{t}$ value is optimal, thereby balancing denoising efficacy and computational efficiency. Additional ablation studies are presented in Appendix to discuss in-depth analysis of the scalability, robustness, and effectiveness of our methods.  
%

%Finally, we assess the performance with varying the negative prompt weight in text-based methods. To establish that our approach consistently enhances performance across diverse setups when integrated with existing methodologies, we conduct a series of experiments. \figref{ablation}‑(c) shows as the weight of SLD increases from Weak to Max, SLD performs inadequately in MMA-Diffusion prompts. In contrast, our safe denoiser improves performance and widens the performance gap, underscoring its robustness and effectiveness as a superior enhancement.
%

\paragraph{Computation Overhead}

\begin{wraptable}{r}{0.35\columnwidth}
  \vskip -0.2in
\caption{Wall-clock time.}
  \label{tab:wall_clock}
\begin{tabular}{lc}
\toprule
\multicolumn{1}{c}{Models}                 & \multicolumn{1}{c}{\begin{tabular}[c]{@{}c@{}}Time \\ (s/img) \end{tabular}} \\
\midrule
SD-v1.4                    & 3.18         \\
\cc{15} + Ours ($N=515$)       & \cc{15}3.20         \\
\midrule
SAFREE                & 4.22         \\
\cc{15} + Ours ($N=515$)   & \cc{15}4.24         \\
\cc{15} + Ours ($N=3,000$) & \cc{15}4.29         \\
\bottomrule
\end{tabular}
\vskip -0.1in
\end{wraptable}
\tabref{wall_clock} presents the wall-clock time for image generation on NVIDIA RTX4090 with 24GB memory. Thanks to GPU parallelism, the additional time introduced by our method scales sub-linearly since modern GPUs optimize batched matrix multiplications with efficient job scheduling. For example, our method increases from 4.22s to 4.29s when using $3,000$ negative images (an overhead of only 0.07s), while increasing from 4.22s to 4.24s when using $515$ images. Conversely, SAFREE adds over 1s per image. Given that SAFREE and ours perform similarly in \tabref{sole}, our method shows a better performance-efficiency curve.

\subsection{Class-Conditional and Unconditional Generation with Safe Denoiser}

\begin{wraptable}{r}{0.35\columnwidth}
  \vskip -0.2in
\caption{Performance in FFHQ. We use ResNet18~\citep{he2016deep} to classify the sex of generated samples. We compute FID by comparing male data and generated images.}
    \label{tab:ffhq}
    \centering
    % \resizebox{0.70\columnwidth}{!}{%
    \resizebox{\linewidth}{!}{
    \begin{tabular}{lccc}
        \toprule
        Models & Female $\downarrow$ & Male $\uparrow$ & FID $\downarrow$ \\\midrule
        Baseline (B)         & $64.0\%$ & $36.0\%$ & 109.07  \\
        B + SR   & $53.1\%$ & $46.9\%$ & 130.52   \\
        \cc{15}B + Ours       & \cc{15}$55.6\%$ & \cc{15}$44.4\%$ & \cc{15}96.57   \\
        \bottomrule
    \end{tabular}
    }
    \vskip -0.1in
\end{wraptable}
This subsection evaluates the performance of safe denoiser when applied in isolation. Specifically, to assess our safe denoiser in a simplified setting, we conduct experiments on two distinct tasks: a class-conditional model trained on ImageNet~\citep{russakovsky2015imagenet} for removing a specific class (e.g., Chihuahua); and on an unconditional model trained on FFHQ~\citep{karras2019style} to negate generating specific sex (e.g., female), see Appendix~\ref{sec:pixel_detail} for details of experiments. For the class removal, \tabref{imagenet_side_by_side} presents precision, recall~\citep{kynkaanniemi2019improved}, and Top-1* (classification accuracy of generated images conditioned by Chihuahua class) metrics for negating the Chihuahua class. As conventional text-based safety techniques are not directly applicable, we compare our method against Sparse Repellency (SR), as described in \secref{related}. \tabref{imagenet_side_by_side} showcases that our method outperforms SR in terms of precision, recall, and Top-1*, indicating that ours avoid generating Chihuahua while being more diverse and precise than SR. 
In the FFHQ experiments, where we targeted the negation of female images, \tabref{ffhq} indicates that while SR exhibits classification results that deceptively suggest successful negation, the FID scores and qualitative comparisons in Appendix~\ref{sec:additional_result} demonstrate that this apparent achievement comes at the cost of significantly degraded image quality. Indeed, our experiments show that our methodology consistently produces visually convincing samples, whereas SR frequently generates out-of-distribution images with artifacts.

% In FFHQ, we aim to prevent the generation of a specific sex. However, since the whole data points are used in training, we select 1K female images from CelebA~\citep{liu2015faceattributes} validation split to serve as unseen negative data, thereby establishing the negative dataset $\{\mathbf{x}^{(1)},...,\mathbf{x}^{(1000)}\}$. We then employ our safe denoiser to sample 1K images. As shown in Table~\ref{tab:ffhq}, classification accuracy for these generated samples reveal that our method more effectively avoids the female class compared to the baseline pretrained model. According to Table~\ref{tab:ffhq}, our algorithm generates more male images compared to Sparse Repellency (SR) while achieving a lower FID score. This suggests that the images produced by SR are of lower quality compared to those generated by ours, leading to increased confusion for the classifier.

% In ImageNet, we focus on negating a specific Chihuahua class during generation. We select the validation set of Chihuahua class as the negative images. We generate 50 samples per class and classify samples from 999 classes by a classifier~\citep{dhariwal2021diffusion} and report the accuracy by Top-1. Also, we measure the Top-1 accuracy of 50 samples from Chihuahua class, reporting it by Top-1* in \tabref{imagenet}. From the result, we note that our method excels generating other 999 classes, while SR cannot generate images from those 999 classes. To evaluate the overall quality, \tabref{imagenet} further report the precision (sample accuracy) and recall (sample diversity)~\citep{kynkaanniemi2019improved} over 50K samples, indicating that our method is better than SR in negating a specific class. 

%

%
%
%
%
%
%
%
%
%
%
%
%
%
%
%
%
%
%
%
%
%



% \subsection{Pixel-Space Generation: Unconditional and Class Conditional Cases}
% \label{sec:pixel}

% In this section, we use our safe denoiser in the DMs without text inputs. Specifically, we employ experiments on FFHQ~\cite{karras2019style} and ImageNet~\cite{russakovsky2015imagenet} in the $256\times256$ resolution. We utilize the pretrained diffusion models from \cite{chung2022diffusion} for FFHQ and \cite{dhariwal2021diffusion} for ImageNet. For the experiments, we use a DPM solver~\cite{lu2022dpm} with 100 steps.

% In FFHQ, we aim to prevent the generation of a specific sex. However, since the whole data points are used in training, we select 1K female images from CelebA~\cite{liu2015faceattributes} validation split to serve as unseen negative data, thereby establishing the negative dataset $\{\mathbf{x}^{(1)},...,\mathbf{x}^{(1000)}\}$. We then employ our safe denoiser to sample 1K images. As shown in Table~\ref{tab:ffhq}, classification accuracy for these generated samples reveal that our method more effectively avoids the female class compared to the baseline pretrained model. According to Table~\ref{tab:ffhq}, our algorithm generates more male images compared to Sparse Repellency (SR) while achieving a lower FID score. This suggests that the images produced by SR are of lower quality compared to those generated by ours, leading to increased confusion for the classifier.

% In ImageNet, we focus on negating a specific Chihuahua class during generation. We select the validation set of Chihuahua class as the negative images. We generate 50 samples per class and classify samples from 999 classes by a classifier~\cite{dhariwal2021diffusion} and report the accuracy by Top-1. Also, we measure the Top-1 accuracy of 50 samples from Chihuahua class, reporting it by Top-1* in \tabref{imagenet}. From the result, we note that our method excels generating other 999 classes, while SR cannot generate images from those 999 classes. To evaluate the overall quality, \tabref{imagenet} further report the precision (sample accuracy) and recall (sample diversity)~\cite{kynkaanniemi2019improved} over 50K samples, indicating that our method is better than SR in negating a specific class. 
%
%

\subsection{Compatibility with Frontier Model and Style-Level Intellectual Property Control}

We evaluate the compatibility of our plug-and-play approach with the frontier model SD-v3 \citep{esser2024scaling}. The experimental results are presented in \tabref{sd3_exp}.
On SD-v3, SAFREE alone reduces ASR by approximately 9\% compared to the baseline. In contrast, our approach achieves a 33.2\% relative reduction in ASR, from 0.304 to 0.203. Notably, our method maintains CLIP alignment and even slightly improves FID. This demonstrates the applicability of our proposed method to recent and powerful backbones models.

Interestingly, our safe denoiser enhances sample diversity during inference, which can lead to a reduction in FID.
A well-known phenomenon of large CFG values is a fidelity and diversity trade-off. Specifically, increasing CFG sharpens alignment but diminishes sample diversity, resulting in a degradation of FID at high values. This phenomenon has been observed in previous studies \citep{sadat2024cads, kynkaanniemi2024applying}.
In contrast, our safe denoiser is not overly reliant on the text conditioning, allowing it to introduce relevant stochasticity that effectively mitigates the loss of diversity caused by high CFG. Consequently, our denoiser improves FID.
Empirically, we have observed higher intra-prompt diversity compared to the baseline. Another perspective to consider is that FID’s Gaussian approximation of feature distributions possibly records small improvements that does not translate into noticeable quality differences in practical applications.

\begin{wraptable}{r}{0.55\columnwidth}
  \vskip -0.2in
    \caption{SD-v3 results on Ring-A-Bell for safety and COCO-30K for image quality.}
    \label{tab:sd3_exp}
    \centering
    % \resizebox{0.7\textwidth}{!}{%
        \begin{tabular}{lcccc}
            \toprule
            \multirow{2}{*}{Method} 
            & \multicolumn{2}{c}{Ring-A-Bell} 
            & \multicolumn{2}{c}{COCO-30K} \\
            \cmidrule(lr){2-3} \cmidrule(lr){4-5}
            & ASR $\downarrow$ & TR $\downarrow$ 
            & FID $\downarrow$ & CLIP $\uparrow$ \\\midrule
            SD-v3      & 0.304 & 0.330 & 23.15 & \textbf{31.46} \\
            + SAFREE & 0.278 & 0.298 & 22.99 & 31.24 \\
            \cc{15}+ Ours   & \cc{15}\textbf{0.203} & \cc{15}\textbf{0.267} & \cc{15}\textbf{22.54} & \cc{15}31.15 \\
            \bottomrule
        \end{tabular}
    % }%
\end{wraptable}

We conceptually evaluate baselines in situations where pretrained diffusion models leak intellectual property. In this scenario, intellectual property-sensitive prompts can be grouped into three scenarios:
\emph{(i) the prompt explicitly names the target intellectural property; 
(ii) the prompt avoids the name but gives a detailed textual description; and
(iii) neither name nor descriptive cues are present, yet the model reproduces the target’s style.}
The third case presents a significant challenge for text-only defenses, as there is no negative text cue to negate. 
As reported by \citep{somepalli2023understanding}, diffusion models can overfit styles and reproduce
them even without explicit textual mentions. We reproduce this phenomenon with Munch’s The
Scream by using the prompt ``If Barbie were the face of the world’s most famous paintings''. While this text prompt never mentions Munch or The Scream, SD-v1.4 recreates the painting’s distinctive
style as shown in \figref{memorized_munch}. When we use four original paintings of ''The Scream'', for instance two from 1893, one from 1895, one
from 1910, as the negative set, our safe denoiser suppresses Munch's style while preserving the
Barbie concept. Additionally, Ours with SAFREE produces both modern and classical renderings without the style of Munch portraits. The qualitative results are displayed in \figref{protected_ip}.

\begin{figure}[!t]
    \centering
    \vskip -0.1in
        \begin{subfigure}{0.49\textwidth}
            \includegraphics[width=\textwidth]{Figures/Experiments/munch.pdf}
            \caption{Negative datapoints}
            \label{fig:munch}
        \end{subfigure} 
        \begin{subfigure}{0.49\textwidth}
                \includegraphics[width=\textwidth]{Figures/Experiments/babie.pdf}
            \caption{Generated Images}
            \label{fig:memorized_munch}
        \end{subfigure}
    \caption{Qualitative result for style-level intellectual property control. SD-v1.4 reproduces Munch’s style, whereas Ours with and without SAFREE removes that style while preserving the ``Barbie'' concept. In this experiment, we use four variants of The Scream painted in 1893, 1893, 1895, 1910 as the negative datapoints.}
    \label{fig:protected_ip}
    \vspace{-0.1in}
\end{figure}


\section{Limitations and Conclusions}
We introduce the \textit{safe denoiser}, an in‑process, training‑free mechanism that steers diffusion model sampling toward theoretically safe distributions, thereby promoting appropriate content. Unlike purely discriminative pre‑ or post‑filters, our approach acts during inference and complements existing guardrails. In particular, this mid‑generation intervention mitigates failure modes in static text or image filters, especially under adversarial prompt engineering. Thus ours contributes to a defense‑in‑depth safety architecture suitable for real‑world applications. Regarding negative datasets, the data requirement is shared across other defenses method. The datasets used to train or calibrate pre‑ and post‑filters can be reused to provide data‑driven negative guidance at inference. A current limitation is the need to tune hyperparameters to balance fidelity and safety. Appendix~\ref{sec:limitation} discusses these trade‑offs and offers practical guidance. Privacy-sensitive generation remains an ongoing challenge, partly due to the lack of standardized quantitative metrics. We leave the development of such metrics and extensions to other modalities for future work. 
% The \textit{safe denoiser} offers a robust and scalable mid-generation safety layer that enhances the security of deployed systems when combined with prompt and output filtering.
\label{sec_conclusion}

\section*{Acknowledgments}
We thank our anonymous reviewers for their constructive feedback, which has helped significantly improve
our paper. We thank the Digital Research Alliance of Canada (Compute Canada) for its computational
resources and services. 
%
M. Kim was supported by the Canada CIFAR AI Safety Catalyst grant.
%
A. Yusuf was funded by the Canada Graduate Scholarships — Master’s program of the Natural Sciences and Engineering Research Council of Canada (NSERC).
%  
M. Park was supported in part by the Natural Sciences and Engineering
Research Council of Canada (NSERC) and the Canada CIFAR AI Chairs program.
% \section*{Impact Statements}

% This paper presents a work whose goal is to build a reliable and trustworthy Generative AI. There are many potential societal consequences of our work, particularly in addressing ethical risks associated with generative models. Our research is focused on preventing the generation of NSFW content, including nudity and violence, and mitigating the risk of models memorizing and reproducing private information, such as human face, from training datasets. We believe the presented work  contributes to the responsible use of generative AI, reinforcing ethical safeguards and promoting AI systems that align with societal values and human rights.

%

%




%



\bibliography{main}
\bibliographystyle{unsrt}

%%%%%%%%%%%% Checklist %%%%%%%%%%%%
\clearpage
\input{Manuscript/checklist}

%%%%%%%%%%%% Appendix %%%%%%%%%%%%
\newpage
\appendix
\setcounter{equation}{0}
\renewcommand{\theequation}{\Alph{section}.\arabic{equation}}
\setcounter{table}{0}
\renewcommand{\thetable}{\Alph{section}.\arabic{table}}
\setcounter{figure}{0}
\renewcommand{\thefigure}{\Alph{section}.\arabic{figure}}

\input{Appendix/Proof}
\section{Limitations and Broader Impacts}
\label{sec:limitation}

\paragraph{Limitations}
We have addressed significant safety challenges in DMs, particularly concerning the generation of NSFW content and the inadvertent reproduction of sensitive data. 
We introduce the \textit{safe denoiser}, a novel approach that modifies the sampling trajectories of DMs to adhere to theoretically safe distributions, thereby ensuring the generation of appropriate and authorized content. 

However, this approach necessitates the introduction of an additional hyperparameter, $\beta_t$, as outlined in \thmref{safe}. 
While we demonstrate that this parameter is theoretically derived and straightforward to implement, it may not be optimal for realistic scenarios due to its assumption of access to numerous data points sampled from an unsafe distribution. 
In practice, we present evidence in \figref{ablation_2} that this parameter influences the performance of the model.

Despite the challenges, we have developed a novel training-free method that effectively guides the sampling trajectories of DMs towards safe distributions. 
Ultimately, this work provides a robust and scalable solution for mitigating safety risks in generative AI, paving a way for their responsible and ethical applications.


\paragraph{Broader Impacts}
This paper presents a work whose goal is to build a reliable and trustworthy Generative AI. 
There are many potential societal consequences of our work, particularly in addressing ethical risks associated with generative models. 
Our research is focused on preventing the generation of NSFW content, including nudity and violence, and mitigating the risk of models memorizing and reproducing private information, such as human face, from training datasets. 
We believe the presented work contributes to the responsible use of generative AI, reinforcing ethical safeguards and promoting AI systems that align with societal values and human rights.
\section{Experimental Details : Text-to-Image Generation}
\label{sec:t2i_detail}

As outlined in the manuscript, we conduct the Text-to-Image experiment using SD-v1.4, following the same model as the baselines for generating images from text, as referenced in \cite{schramowski2023safe, wu2024erasediff, gong2024reliable, yoon2024safree}. 
To ensure consistency, we adopt the generation procedure described in each baseline. 
Preliminary observing the sensitivity of nudity-related content, we employ the DDPM scheduler \cite{ho2020denoising}. %
For a fair comparison, we maintain the same number of inference steps, specifically $50$, aligning with the official implementations of both SLD and SAFREE, which also use 50 inference steps.

Regarding the \textit{Safe Denoiser}, the proposed model computes the transition kernel with an RBF kernel. 
%
%
The RBF kernel function is defined as follows:
\begin{equation}
\label{eq_rbf_kernel}
K(x, x') = \exp\left(-\frac{\lVert x - x'\rVert^2}{2\sigma^2}\right)
\end{equation}
For the bandwidth parameter $\sigma$, we set a value of 1.0 for SLD and 3.15 for SAFREE. 
Additionally, in case of SAFREE, we apply a scaling factor $\eta=0.33$, whereas for SLD, we use $\eta=0.03$ to regulate the strength of the repellency in \eqref{uncond_safe}.
Empirically, we introduce a heuristic in which the proposed \textit{Safe Denoiser} is applied within critical timesteps $C=[780, ...,1000]$. %
In the early stages of diffusion, denoising process primarily establishes global structures rather than intricate details, while the later stages focus on refining fine-grained features. 
Since our approach aims to prevent the generation of globally harmful images rather than enhancing image quality or detail, we apply the denoiser at these later timesteps. 

For reference images, we provide a detailed explanation of how they are obtained. 
To ensure safe generation against nudity prompts, we utilize a total of 515 images sourced from the I2P dataset~\cite{schramowski2023safe}. 
These images were generated using SD-v1.4. As mentioned in the manuscript, these reference images meet the criterion, where a nude class probability exceeds 0.6, as determined by Nudenet. 
Sample images are presented in \figref{i2p_ref_imgs_1}. 
On the other hand, for the inappropriate probability task with the CoPro dataset, we attempt to use the total images from the I2P dataset. However, our computational resources allow us to use only 3,000 reference images.
To select these 3,000 images, we randomly choose them out of the 4,703 images available in the I2P dataset. All images used in this task are also generated using SD-v1.4.
Sample images are presented in \figref{i2p_ref_imgs_2}. 
It’s important to note that all experiments conducted in this study use the same set of reference images across all baselines. This ensures a fair comparison.

% Please delete the following parts for arXiv version
Addtionally, the reference images we used in the task to generate safer images against nudity prompts are included as attachments in our supplementary materials.
On the other hand, due to space constraints, we cannot include the attachment in the inappropriate probability task. We ensure that the reference images used in this task will be included in the public repository upon the acceptance of the paper.

\begin{figure}[h]
    \centering
    \includegraphics[width=0.92\textwidth]{Figures/Appendix/Implementation_detail/ref_imgs.pdf}
    \caption{Reference images for safe generation against nudity prompts}
    \label{fig:i2p_ref_imgs_1}
\end{figure}

\begin{figure}[h]
    \centering
    \includegraphics[width=0.92\textwidth]{Figures/Appendix/Implementation_detail/ref_imgs_whole_i2p.pdf}
    \caption{Reference images for inappropriate propability}
    \label{fig:i2p_ref_imgs_2}
\end{figure}

Next, we briefly introduce the baseline models used in our experiments.
The first two approaches serve as comparisons for unlearning-based safe diffusion models \cite{gandikota2023erasing, gong2024reliable}.  
Specifically, we evaluate Erased Stable Diffusion (ESD) \cite{gandikota2023erasing} as a representative method. More recently, reliably trained safe diffusion (RECE) models have demonstrated improved performance, particularly in reducing the attack success rate \cite{gong2024reliable}. 
In addition to these unlearning-based approaches, we also include SLD and SAFREE as training-free safe diffusion models \cite{schramowski2023safe, yoon2024safree}. While both methods employ negative prompts, their underlying mechanisms differ significantly. 
In SLD, the set of unsafe prompts, denoted as  $c_{US}$, is designed to mitigate globally harmful image generation \cite{schramowski2023safe}.
In contrast, SAFREE focuses on more precise negative prompts specifically tailored to nudity-related content \cite{yoon2024safree}. Beyond negative prompts, SAFREE further enhances safety by applying an orthogonal projection technique in Euclidean space to shift text embeddings away from predefined toxic regions.
In the following, we provide an overview of the datasets used in our experiments.

\subsection{Inappropriate Prompt Datasets}

\textbf{I2P}  
The I2P dataset consists of prompts related to seven unsafe concepts: hate, harassment, violence, self-harm, sexual content, shocking content, and illegal activity \cite{schramowski2023safe}. 
It contains a total of 4,703 prompts and was introduced in earlier stages of research, with subsequent studies primarily focusing on this dataset \cite{gong2024reliable, yoon2024safree}. 
In this work, we utilize the I2P dataset as a source of reference data points rather than for additional training.
The dataset was obtained from \url{https://huggingface.co/datasets/AIML-TUDA/i2p}

\textbf{CoPro}
Compared to I2P \cite{schramowski2023safe}, the CoPro dataset offers a more extensive dataset comprising a total of 226,104 prompts, each associated with 723 concepts that span both safe and unsafe scenarios. This expansion enhances the dataset’s suitability for rigorous evaluation \cite{liu2024latent}. Particularly, it also offers super-concept information, following the same framework of I2P \cite{schramowski2023safe}. All text prompts are categorized into \{hate, harassment, violence, self-harm, sexual content, shocking content, and illegal activities\}. This ensures that they align with the corresponding category information in the I2P dataset.
To efficiently evaluate models, we randomly sample 10,000 prompts, ensuring a uniform distribution across all categories. We validated that the average inappropriate probability of SD-v1.4 in the randomly sampled dataset, presented in \tabref{ip_CoPro}, closely matches the numerical information provided in \cite{liu2024safetydpo}.
In this work, we evaluate safe image generation performance across baselines on the CoPro dataset using reference data points from the I2P dataset.
This dataset was obtained from \url{https://github.com/rt219/LatentGuard/blob/main/dataset/CoPro_v1.0.json}

\subsection{Nudity in NSFW Prompt Datasets}

\textbf{Ring-A-Bell}
The Ring-A-Bell dataset was developed through a red-teaming approach that evaluates text-to-image diffusion models using black-box methods \cite{tsai2024ringabell}. 
The original dataset \url{Chia15/RingABell-Nudity} contains 285 prompts; however, we use a curated subset of 79 prompts, following prior baselines \cite{gong2024reliable, yoon2024safree}. 
This selection ensures a more equitable comparison of our method. 
The curated Ring-A-Bell dataset was obtained from either \url{https://github.com/CharlesGong12/RECE} or \url{https://github.com/jaehong31/SAFREE}.

\textbf{MMA-Diffusion}
MMA-Diffusion is another dataset generated via a red-teaming approach \cite{yang2024mma}. 
Unlike other datasets, it consists of adversarial prompts designed to include potentially harmful contexts without explicit expressions. 
Similar to the Ring-A-Bell dataset, we use a curated set of 1,000 prompts, consistent with prior baselines \cite{gong2024reliable, yoon2024safree}. 
The dataset was obtained from \url{https://github.com/CharlesGong12/RECE} or \url{https://github.com/jaehong31/SAFREE}.

\textbf{UnlearnDiff}
The UnlearnDiff dataset contains various harmful text prompts that can potentially generate NSFW images \cite{zhang2024generate}. 
Among its categories, we specifically focus on nudity-related prompts. 
The dataset includes a total of 116 nudity-related prompts, derived from an initial set of 143 prompts, from which 27 were excluded as they contained other NSFW categories such as self-harm and shocking content. 
This selection ensures that our numerical metrics remain unaffected by unrelated factors.
The dataset was obtained from \url{https://github.com/CharlesGong12/RECE} or \url{https://github.com/jaehong31/SAFREE}.

\subsection{Ann Graham Lotz for Data Memorization}

In \figref{thumbnail_2}, we demonstrate that SD-v1.4 exhibits trainig dataset memorization, as it is capable of regenerating an indentical images using the text prompt, (\textit{'Living in the light with Ann Graham Lotz <|startoftext|> lad mans'}). In this example, our method is applied with a bandwidth $\sigma=13.15$ and scaling factor of $0.69$. 
To construct a reference data for this case, we collected a total of 10 images from the internet. These are presented in \figref{ann_ref_imgs}
\begin{figure}[h]
    \centering
    \includegraphics[width=0.60\textwidth]{Figures/Appendix/Implementation_detail/ref_imgs_Ann.pdf}
    \caption{Reference images for Ann Graham Lotz case}
    \label{fig:ann_ref_imgs}
\end{figure}

\section{Experimental Details : Class-Conditional and Unconditional Generation}
\label{sec:pixel_detail}

In this section, we use our safe denoiser in the DMs without text inputs. 
Specifically, we employ experiments on FFHQ~\cite{karras2019style} and ImageNet~\cite{russakovsky2015imagenet} in the $256\times256$ resolution. 
We utilize the pretrained diffusion models from FFHQ~\cite{chung2022diffusion}\footnote{https://github.com/DPS2022/diffusion-posterior-sampling} and ImageNet~\cite{dhariwal2021diffusion}\footnote{https://github.com/openai/guided-diffusion}. 
For the experiments, we use a DDPM solver~\cite{lu2022dpm} with 100 steps.

\paragraph{Unconditional Generation}
For unconditional generation, we utilize the FFHQ dataset to evaluate whether the proposed method effectively mitigates sexual bias, using our method. 
Although FFHQ datset does not include explicit label information, \tabref{ffhq} illustrates that the generated images exibit a noticiable bias toward female images over male ones. 
In this experiment, we select 1K female images from CelebA-HQ~\citep{liu2015faceattributes}\footnote{https://www.kaggle.com/datasets/badasstechie/celebahq-resized-256x256} validation split to serve as unseen negative data, thereby establishing the negative dataset $\{\mathbf{x}^{(1)},...,\mathbf{x}^{(1000)}\}$. 
Then, we employ our safe denoiser to generate 1K images.
While both FFHQ and CelebA-HQ are designed to capture similar distribution, they are not completely aligned. This distinction provides an advantageous experimental setup, where we assess the controllability of image generation using reference images. 
For performance evaluation, we compute FID~\cite{heusel2017gans} score using 1,000 male images from the CelebA-HQ dataset. 
For classification performance, we train a ResNet18 model, as implemented in the PyTorch framework~\footnote{\url{https://pytorch.org/vision/stable/index.html}} using the training dataset in the CelebA-HQ. 
In this experiment, we chose $\sigma=1.0$ and $\eta=0.05$, and employ \textit{Safe denoiser} across the entire denoising timesteps. 

\paragraph{Conditional Generation}
For conditional ImageNet~\cite{russakovsky2015imagenet} experiments at $256 \times 256$ resolution, we use a diffusion model trained on the full ImageNet-256 dataset guided by a classifier~\cite{dhariwal2021diffusion}. 
The diffusion backbone follows a linear noise schedule and is constructed with 1,000 diffusion time-steps. We condition on class labels by scaling the classifier guidance at 5.0, creating a strong pull towards the desired class during the sampling process. 
Each experiment generates 50 samples per class across all 1000 ImageNet classes, producing 50,000 samples that are then evaluated with a pretrained ImageNet classifier for precision, recall, and classification accuracy measurements~\cite{he2016dppresnet50}. 
Our metrics include \textbf{(i) Precision: } the fraction of generated samples that match the designated ImageNet label when conditioned on the class, \textbf{(ii) Recall: } aims to evaluate the diversity and coverage of the targeted class distribution, and \textbf{(iii) Classification Accuracy: } the rate at which generated images are correctly identified as their conditioned label among the 999 classes (excluding the negated target class, i.e, Chihuahua). 
The classification accuracy on the hold-out negated class is also calculated, to evaluate how well the respective method does not generate the negated target class. 
As illustrated in \tabref{imagenet_side_by_side}, we focus on the Chihuahua class to investigate how effectively our proposed safe denoiser can repel a target class while preserving generative quality for other classes in this experiment. 
To avoid unintended Chihuahua generation, aforementioned metrics aim to make sure that samples do not drift toward distinct Chihuahua-like features.
For instance, when we generate an image based on a reference dataset sampled from a Chihuahua, the resulting sample may resemble a Golden Retriever, but it won’t resemble a Chihuahua.

To compare our approach, we implement three variants of the conditional diffusion process: vanilla classifier-guided diffusion model without repellency mechanisms, the Sparse Repellency (SR)~\cite{kirchhof2024sparse} technique applied to the classifier-guided diffusion model, and our safe denoiser technique applied to the same diffusion process. 
For the reference dataset, we select the validation set of Chihuahua class as the negative images.
In this experiment, the safe denoiser technique is applied on the 200 to 800 timesteps of the diffusion process. $\eta=0.02$ was chosen as to control the strength of the repellency away from the Chihuahua target class. 
In the SR variant of the experiment, a repellency scale of 0.01 is combined with a large radius of 300 to push generated samples out of regions resembling the negated target class.
\section{Additional Experimental Results}
\label{sec:additional_result}

In this section, we share extra experimental results. Both numeric and visual results are included, which are not presented in the main text. These results highlight the empirical gains in terms of safe generation and the preservation of the global context of the generated samples simultaneously. Specifically, this ensures that the samples remain faithful to their original meanings while enabling us to negate specific concepts we intended. 

\subsection{Quantiative Results}

\paragraph{ImageNet Case for Negating Chihuahua Class} 
In ImageNet, we focus on negating a specific Chihuahua class during generation. We select the validation set of Chihuahua class as the negative images. We generate 50 samples per class and classify samples from 999 classes by a classifier~\citep{dhariwal2021diffusion} and report the accuracy by Top-1. Also, we measure the Top-1 accuracy of 50 samples from Chihuahua class, reporting it by Top-1* in \tabref{imagenet}. From the result, we note that our method excels generating other 999 classes, while SR cannot generate images from those 999 classes. To evaluate the overall quality, \tabref{imagenet} further report the precision (sample accuracy) and recall (sample diversity)~\citep{kynkaanniemi2019improved} over 50K samples, indicating that our method is better than SR in negating a specific class. 

\begin{table}[!ht]
    \caption{Experiments on ImageNet for the specific class (Chihuahua) negation task. Top-1 is the classification accuracy of the generated samples on 999 classes, and Top-1* indicates the accuarcy on the specific class.} 
    \label{tab:imagenet} 
    \centering 
    \begin{tabular}{lccccc} 
        \toprule 
        Method & Prec $\uparrow$ & Rec $\uparrow$ & Top-1 $\uparrow$ & Top-1$^{*}$ $\downarrow$ \\ 
        \midrule 
        Baseline (B) & 0.72 & 0.63 & 0.76 & 0.68 \\ 
        B + SR & 0.59 & 0.54 & 0.01 & 0.0 \\ 
        \cc{15}B + Ours & \cc{15}0.62 & \cc{15}0.58 & \cc{15}0.14 & \cc{15}0.0 \\ \bottomrule 
    \end{tabular} 
\end{table}

\paragraph{Aesthetic Scores for Long Text Prompts} 
We identified that our method, which incorporates negative prompts, effectively reduces the risk of generating unsafe data and maintains alignment with the given text prompts. However, the text prompts in these cases span various lengths. Therefore, it is necessary to quantify whether our method excessively applies to remove unsafe contents, leading to unfavorable images in extreme cases. We sample the most complex cases from the I2P datasets and compare the generated images across baselines. 

\begin{table}[!ht]
    \caption{Aesthetic scroes for long text prompts in the I2P dataset.} 
    \label{tab:aes} 
    \centering 
    \begin{tabular}{lc} 
    \toprule 
    Method      & LAION-aesthetic V2 $\uparrow$ \\
    \midrule 
    SD-v1.4          & $5.97 \pm 0.534$         \\
    SAFREE      & $6.03 \pm 0.540$              \\
    SAFREE+Ours & $5.94 \pm 0.529$              \\
    \bottomrule 
    \end{tabular}
\end{table}


To test how our method works when long and complex prompts are given, we use LAION-aesthetic V2 score~\footnote{https://github.com/christophschuhmann/improved-aesthetic-predictor} as a metric and use top 10\% longest prompts (475 prompts, avg. word\_count=54) selected from the I2P dataset. We choose this score as it is known to be correlated with human perception of quality of images (higher the better). As shown \tabref{aes}, our method maintains aesthetic quality comparable to baselines, even with complex prompts. We prove that the proposed method does not struggle to create appropriate images even when asked with long text prompts.  

\subsection{Qualitative Results}
We present additional qualitative results across three experimental scenarios: \textit{(1) Text-to-Image Generation for preventing nudity and inappropriate images, (2) Sexual Debiasing in unconditional generation for facial images, and (3) Class-Conditional Generation, where reference images serve as constraints not to generate.} To systematically demonstrate the effectiveness of our approach, we present the results in sequence, beginning with text-to-image generation followed by unconditional generation and concluding with conditional generation. To facilitate straightforward understanding, we include as many figures and qualitative comparisons as possible.

\subsection{Text-to-Image Generation}
\paragraph{Safe Generation against Nudity Prompts}
We present a qualitative comparison across baselines and ours. All figures are generated using the same text prompts. We decide to exclude the case of MMA-Diffusion since prompts in this dataset generate pornographic images by baselines, which is not suitable for academic purpose. Hence, we select text prompts from Ring-A-Bell~\cite{tsai2024ringabell} and UnlearnDiff~\cite{zhang2024generate}. From \figref{ringabell} to \figref{unlearndiff}, we observe that our model effectively eliminates nudity information while preserving textual information. 

\begin{figure*}[!ht]
    \centering
    \includegraphics[width=0.80\textwidth]{Figures/Appendix/Implementation_detail/ref_imgs_ring.pdf}
    \caption{Generated images by baselines and ours on Ring-A-Bell~\cite{tsai2024ringabell}}
    \label{fig:ringabell}
\end{figure*}

\begin{figure*}[!ht]
    \centering
    \includegraphics[width=0.80\textwidth]{Figures/Appendix/Implementation_detail/ref_imgs_unlearn.pdf}
    \caption{Generated images by baselines and ours on UnlearnDiff~\cite{zhang2024generate}}
    \label{fig:unlearndiff}
\end{figure*}


\paragraph{Inappropriate Probability in CoPro Dataset}

In our evaluation on the CoPro dataset~\cite{liu2024latent}, we apply our method with images from the I2P dataset~\cite{schramowski2023safe}, which includes a wide range of sensitive categories: \{hate, harassment, violence, self-harm, sexual content, shocking content, and illegal activities\}. Among these, we focus on the 'Self-Harm' category.
Self-harm content is suitable for graphical illustration, distinguishing it an appropriate and interpretable case for homogenous visual inspection in the public domain. Unlike other categories—where perceptions of appropriateness can vary widely across cultural and personal contexts—'Self-harm' is typically associated with somber or distressing imagery that is broadly and publicly recognized as unsafe.

\begin{figure*}[!ht]
    \centering
    \includegraphics[width=0.75\textwidth]{Figures/Appendix/Implementation_detail/ref_imgs_copro_selfharm.pdf}
    \caption{Generated images by baselines and ours on CoPro~\cite{liu2024latent}. All texual prompts are labeled as 'Self-Harm'. This dataset also provides safe alternatives, and we present both.}
    \label{fig:copro_selfharm}
\end{figure*}

As illustrated in \figref{copro_selfharm}, our Safe Denoiser effectively reduces the generation of implicit unsafe content while preserving coherence with the provided prompts. This underscores its ability to both detect and suppress sensitive content without compromising the semantic alignment of textual prompts.
From this figure, it is evident that negative prompts do not yield significant results in mitigating the generation of sad and gloomy atmospheres, particularly for conveying feelings of collapse. Conversely, our \textit{Safe Denoiser} tends to generate images that more accurately reflect the literal meanings of the textual prompts. This tendency contributes to a reduction in the likelihood of generating content that evokes feelings of 'Self-harm'. 


\paragraph{Alignment of Textual Prompts}
We present uncurated generated images from the CoCo dataset. This dataset encompasses a wide range of textual prompts that cover various lengthy and diverse situations. As shown in \figref{coco30k_ours} and \tabref{t2i}, we conclude that our approach does not compromise the performance of generating normal images. Instead, it focuses on addressing the challenge of generating potentially unsafe images. 

\begin{figure*}[!ht]
    \centering
    \includegraphics[width=0.70\textwidth]{Figures/Appendix/Implementation_detail/ref_imgs_ours_coco30k.pdf}
    \caption{Uncurated generated images by SAFREE+Ours on CoCo30K}
    \label{fig:coco30k_ours}
\end{figure*}

\subsection{Unconditional Generation: Sexual Debiasing}
We present uncurated generated images created by the DM trained on the FFHQ dataset~\cite{karras2019style}. This dataset lacks explicit labels indicating sexual information. However, we observe a tendency for this model to generate female images more frequently than male images, as shown in \tabref{ffhq}. On the other hand, when we utilize \textit{Safe denoiser} with female images, we mitigate the potential bias towards female images and achieves generating images uniformly distributed across sexual information. \figref{ffhq_comparision} illustrates that the generated images align with the numerical results presented in \tabref{ffhq}.

\begin{figure*}[!ht]
    \centering
    \begin{subfigure}{0.32\textwidth}
            \includegraphics[width=\textwidth]{Figures/Appendix/Implementation_detail/ref_imgs_uncond.pdf}
            \caption{Uncondtional FFHQ}
    \end{subfigure}   
    \begin{subfigure}{0.32\textwidth}
            \includegraphics[width=\textwidth]{Figures/Appendix/Implementation_detail/ref_imgs_uncond_sr.pdf}
            \caption{Sparse Repellency}
    \end{subfigure}   
    \begin{subfigure}{0.32\textwidth}
            \includegraphics[width=\textwidth]{Figures/Appendix/Implementation_detail/ref_imgs_uncond_ours_1.pdf}
            \caption{Ours}
    \end{subfigure}   
    \caption{Comparison of \textit{Safe Denoiser} against existing approaches when negation on female.}
    \label{fig:ffhq_comparision}
\end{figure*}



\subsection{Class-conditional Generation : Negation of Specific Class}
We present uncurated images that focus on negating a specific Chihuahua class. Here are two experimental setups. First, we employ class conditional guidance on the ‘Chihuahua’ class and simultaneously use the \textit{Safe denoiser} to work with negative images sampled from the ‘Chihuahua’ class in the validation split. We observe that the SR does not follow homogeneous images that align with the superclass, ‘Dog’, but our method produces similar small dogs but not matched with ‘Chihuahua’ as shown in \figref{imagenet_comparison_dogs}.

Second, we qualitatively evaluate that our method with negative images from the ‘Chihuahua’ class works when class guidance is applied to classes other than ‘Dogs’, for example, ‘Tench’ and ‘Truck’. This result is shown in \figref{imagenet_comparison_others}. We observe that the SR sometimes depicts different classes even when class guidance is applied, but our method aligns with homogeneous classes following class guidance even when the \textit{Safe denoiser} works with ‘Chihuahua’ images. This indicates that our method effectively tackles specific concepts and preserves the original performance when it is not mutually correlated. 

\begin{figure*}[!ht]
    \centering
    \includegraphics[width=0.99\textwidth]{Figures/Experiments/imagenet/vs_dogs/comparison_imagenet_chihuahua_vs_dogs.pdf}
    \caption{Generated samples when negating the Chihuahua class, primarily producing visually similar small dog breeds.}
    \label{fig:imagenet_comparison_dogs}
\end{figure*}

\clearpage

\begin{figure*}[!ht]
    \centering
    \includegraphics[width=0.99\textwidth]{Figures/Experiments/imagenet/vs_others/comparison_imagenet_chihuahua_vs_others.pdf}
    \caption{Comparison of \textit{Safe Denoiser} against existing approaches when negation on Chihuahua. This comparison includes non-dog related ImageNet classes, which include Tench, Garbage Truck, Church, Spoonbill, and Great White Shark.}
    \label{fig:imagenet_comparison_others}
\end{figure*}


Additional graphical illustrations are presented in the following figures from \figref{imagenet_vanilla_samples} to \figref{imagenet_safe_denoiser_samples}. 

\begin{figure*}[!ht]
    \centering
    \includegraphics[width=0.92\textwidth]{Figures/Experiments/imagenet/vs_others/vanilla_imagenet_chihuahua_vs_others.pdf}
    \caption{Classifier guidance diffusion model generated samples when negating on Chihuahua. This comparison includes non-dog-related ImageNet classes mentioned in~\ref{fig:imagenet_comparison_others} along with the dog-related classes in Figure~~\ref{fig:imagenet_comparison_dogs} which are Pomeranian, Yorkshire Terrier, and Shih Tzu.}
    \label{fig:imagenet_vanilla_samples}
\end{figure*}

\begin{figure*}[!ht]
    \centering
    \includegraphics[width=0.92\textwidth]{Figures/Experiments/imagenet/vs_others/sparse_repellency_imagenet_chihuahua_vs_others.pdf}
    \caption{\textit{Sparse Repellency} generated samples when negating on Chihuahua. The same classes are selected as ~\ref{fig:imagenet_vanilla_samples}.}
    \label{fig:imagenet_sparse_repellency_samples}
\end{figure*}


\begin{figure*}[!ht]
    \centering
    \includegraphics[width=0.92\textwidth]{Figures/Experiments/imagenet/vs_others/safe_denoiser_imagenet_chihuahua_vs_others.pdf}
    \caption{\textit{Safe Denoiser} generated samples when negating on Chihuahua. The same classes are selected as ~\ref{fig:imagenet_vanilla_samples}.}
    \label{fig:imagenet_safe_denoiser_samples}
\end{figure*}


\end{document}